\documentclass[11pt]{amsart}
\usepackage[T1]{fontenc}
\usepackage{lmodern,amsmath,amssymb,mathtools,microtype,booktabs}
\usepackage[margin=1.15in]{geometry}
\usepackage[hidelinks]{hyperref}
\newtheorem{theorem}{Theorem}[section]
\newtheorem{proposition}[theorem]{Proposition}
\newtheorem{lemma}[theorem]{Lemma}
\newtheorem{corollary}[theorem]{Corollary}
\theoremstyle{definition}\newtheorem{definition}[theorem]{Definition}
\theoremstyle{remark}\newtheorem{remark}[theorem]{Remark}
\newcommand{\R}{\mathbb R}
\newcommand{\Z}{\mathbb Z}
\newcommand{\E}{\mathbb E}
\newcommand{\Prob}{\mathbb P}
\newcommand{\one}{\mathbf 1}
\newcommand{\dd}{\,\mathrm d}
\DeclareMathOperator{\dist}{dist}
\DeclareMathOperator{\covol}{covol}
\DeclareMathOperator{\area}{area}
\DeclareMathOperator{\diam}{diam}
\numberwithin{equation}{section}
\title[Scalar collision laws]{Scalar collision laws and recognition of periodic dispersing billiards}
\author{Qian Qi}
\date{October 4, 2026}
\subjclass[2020]{37D50, 52A10, 60J10, 62G05, 68Q32}
\keywords{Dispersing billiards, reciprocal collision commands, active boundary estimation,
query complexity, unknown launch footprint, stationary noise, scale identification}
\hypersetup{pdftitle={Scalar collision laws and recognition of periodic dispersing billiards},pdfauthor={Qian Qi}}
\begin{document}
\raggedbottom
\begin{abstract}
We recover periodic dispersing tables from attempted collision bits
under a fixed unknown stationary launch law. Two homothetic launch
footprints suffice to determine both the obstacles and the operating
footprint even when the homothety ratio is unknown. An integrated pooled collision mean measures an obstacle width sum,
which determines the ratio by a linear support identity. An independent
area normalization, obtained by a finite killed-walk adjoint, proves
identifiability under a weaker separation condition.
For bounded $C^{6,\beta}$ classes, $s=6+\beta$, and boundary density
bounded below by a constant times distance to the power $\gamma$, we
give a finite construction with attempted-bit upper bound
$C\nu^{-((\gamma+3/2)s+1)/(s-2)}$
up to logarithmic factors at $C^2$ accuracy $\nu$.
The improvement comes from a finite conjunction of translated pooled
collision tests, each retaining the unobserved compass direction.
The tests separate a zero probability from a rare positive probability
without subtracting noisy means. For a fixed uniform-disk launch law,
the retained expected-attempt lower power is $(s+1)/(s-2)$; the
remaining exponent gap is $s/(2(s-2))$.
The construction also bounds command descriptions and nominal precision.
Exact homothety about a fixed laboratory origin and quantitative
geometric and boundary-mass priors are assumed; neither the density nor
a footprint support function is supplied. All solid starts and misses
are charged. Uniform finite primitive-period decisions use a known
positive nonperiod-patch margin. The localized inverse, its matching
binary powers and the bounded-error calibration converse are retained
with complete proofs.
\end{abstract}
\maketitle
\section{The collision experiment}
\label{sec:setting}
We study the recovery of a periodic dispersing table when an attempted
preparation returns only its collision bit. The actual launch position
is not observed, and a solid start has the same output as a free miss.
At each of two settings the launch law is stationary with a fixed,
unknown convex support. The supports are homothetic, but their ratio
need not be supplied. Our principal results reconstruct the table, the
operating footprint and this ratio. The proofs use two features of the
collision experiment: translated pooled commands give one-sided tests
near an exposed boundary, and a reciprocal occupation integral recovers
the area of an individual obstacle.

The experiment and the geometric loss are specified first. A summary
of the stationary results follows these definitions. The localized
experiment and its complete proofs are retained in the appendices, where
they also supply the coarse acquisition and period-recognition tools.

\subsection{Physical class and loss}
Fix $0<\beta\le1$ and put $s=6+\beta$. A table is the nonempty union
\begin{equation}\label{eq:presentation}
 \mathcal O=\bigcup_{i=1}^r(C_i+L\Z^2),\qquad 1\le r\le r_0.
\end{equation}
This presentation has exactly $r$ component orbits. It need not be
primitive. The constants in the following bounds are fixed and known:
\begin{equation}\label{eq:priors}
 \|L\|,\|L^{-1}\|\le L_0,\quad |\det L|\ge V_0>0,\quad
 \diam C_i\le D_0,\quad \dist(C,D)\ge d_0>0
\end{equation}
for distinct physical components. All components are strictly convex,
with curvature in $[\kappa_-,\kappa_+]\subset(0,\infty)$ and centered
support functions bounded by $M$ in $C^{6,\beta}(\R/2\pi\Z)$.
Reflections, repeated shapes and nonprimitive presentations are allowed.
Write $\mathcal B$ for this class, $\Pi=\{w:\mathcal O+w=\mathcal O\}$
for its full period group, $z_C$ for a component's Steiner point, and
$p_C$ for its centered support function. All norms on finite-dimensional
spaces are fixed once and for all.

A known positive nonperiod-patch margin $\eta$ will be required only
for uniform finite discrete decisions. Its precise definition is
\eqref{eq:patch-margin}. We denote the resulting subclass by
$\mathcal B_\eta$. The geometric loss is the maximum $C^2$ error of
matched component support functions in a protected laboratory patch,
together with the error of a matched primitive period basis and free
area. Representatives are chosen in a common protected patch. The
basis is matched to a true basis, not required to be a continuously
chosen reduced basis. Exact rational relations and the primitive orbit
count are stated separately from this real-valued loss.

\subsection{Commands and their accounting}
For a nominal displacement $a$, let $B_a(q)$ equal one if a free start
at $q$ has a first collision on the unreflected segment $[q,q+a]$.
It equals zero for a solid start and for a free miss. Both are attempted
preparations and enter the denominator. For pointwise statements solids
and swept sets are closed: a boundary solid start is zero, and a start
on a swept boundary but outside every solid is one. This fixes the
tangency convention; changing it on null sets has no effect on the means.
Fix an even smooth probability kernel $\kappa$ supported in the unit disk,
and set
\[
 f_{x,\sigma}(q)=\sigma^{-2}\kappa((q-x)/\sigma),\qquad
 u_\sigma=\kappa_\sigma*\one_\mathcal O.
\]
Choose $\sigma_0,t>0$ with $t+2\sigma_0<d_0$, and prescribe the compass
law
\begin{equation}\label{eq:compass}
 \rho=\tfrac14(\delta_{te_1}+\delta_{-te_1}
                         +\delta_{te_2}+\delta_{-te_2}).
\end{equation}
A forward command draws $q\sim f_{x,\sigma}$ and $a\sim\rho$
independently and attempts $(q,a)$. A reverse command draws that same
nominal joint law and attempts $(q+a,-a).$ Separate fresh preparations
are used; pairing the same physical particle is unnecessary. Only the
two pooled means are used, not means resolved by direction. The joint
reverse law, rather than only its marginal direction law, is prescribed.

In this localized bounded-error experiment, each actual preparation may be coupled to its nominal one with position,
duration and angle errors bounded by $\ell,\tau,\alpha$, respectively.
It remains a unit-speed straight flight until its first impact. The
bounds hold conditionally on past commands, and preparations at a given
command are fresh and conditionally independent. Section~\ref{sec:queries}
proves the bias bound $C(\ell+\tau+t\alpha)/\sigma$. Calibration is a
resource of the controller, not an estimate inferred from the bits.
The launch aperture is fixed by the prior and does not grow as accuracy
increases. It includes the translated reverse starts and the kernel
supports. No free point, obstacle label, impact location, impact angle,
collision time, primitive cell or area normalization is supplied.

\subsection{Stationary reconstruction and unknown scale}
In the stationary experiment, a forward attempt at nominal center $x$
starts at $x+Z$, with an independent compass displacement, and a reverse
attempt starts at $x+Z+a$ with displacement $-a$. The draw $Z$ has an
unknown density on an unknown convex body $K$. The second setting has
footprint $rK$, where the unknown ratio satisfies known bounds
$1+g_*\le r\le R_*$. We take the first operating footprint as the
reference unit of dilation. Both supports have fixed positive size;
they do not contract as accuracy increases. Homothety is about the
specified laboratory origin. The numerical inball, envelope, curvature
and $C^{6,\beta}$ bounds and a lower boundary-mass condition are given.
In particular, for fixed $b_0>0$ and $\gamma\ge0$ the first density obeys
$j(z)\ge b_0\dist(z,\partial K)^\gamma$ almost everywhere in $K$.
Section~\ref{sec:unknown-footprint} gives the detailed geometric priors.

Choose once a dyadic compass step so that the largest footprint
diameter plus $2t$ is less than $d_0$. This stricter design leaves a
positive collar for the translated commands used below. It is a
prior-dependent choice of the same four-direction pooled apparatus.
No direction-resolved response, launch position, occupancy bit or
obstacle area is observed.

The principal constructive conclusion is
Theorem~\ref{thm:unk-scale-finite}. It recovers $K$, $r$ and the original
table to laboratory $C^2$ error $C\nu$, and, on $\mathcal B_\eta$, the
primitive discrete data with probability at least $1-\delta$. Put
\begin{equation}\label{eq:new-upper-exponent}
 A_s(\gamma)=\frac{(\gamma+3/2)s+1}{s-2},\qquad
 L_\nu=\log(C/\nu),\qquad L_{\nu,\delta}=\log(C/(\nu\delta)).
\end{equation}
The sufficient attempted-bit and command-description bounds are
\begin{align}
 N_\nu&\le C\nu^{-A_s(\gamma)}L_\nu L_{\nu,\delta},
                                              \label{eq:main-new-attempts}\\
 \mathcal D_\nu&\le C\bigl\{\nu^{-1/(s-2)}L_\nu
                      +\nu^{-2}L_{\nu,\delta}\bigr\}L_{\nu,\delta}.
                                              \label{eq:main-new-control}
\end{align}
Here $\mathcal D_\nu$ is the binary length of the commanded centers,
setting indices and repetition counts. The finest dyadic nominal mesh
is $O(\nu^{s/(s-2)})$. These statements price the input descriptions
alongside the observations; physical manufacture, travel, homothety
certification and arithmetic running time remain separate resources.
All attempted preparations, including solid starts and misses, are
included in $N_\nu$. The numerical priors are known with certified
enclosures; neither $r$ nor a support-function evaluation oracle for
$K$ is supplied.

There are two new steps. First, near each recovered coarse boundary,
a finite conjunction of shifted pooled collision tests has a zero
probability outside and a uniformly positive probability at depth $e$.
Proposition~\ref{prop:rare-query} proves that its cost is
$O(e^{-(\gamma+3/2)}\log(1/\varepsilon))$. This avoids the squared
inverse-signal cost of the signed occupation query. Second, the integral of a pooled forward mean measures the sum of two
obstacle widths. A linear identity with the two positivity supports
then determines the unknown homothety ratio. Its $O(\nu^{-2})$
sampling cost is absorbed
by \eqref{eq:main-new-attempts}, since $6<s\le7$ and
$A_s(\gamma)>2$. As a second exact argument, a finite killed-walk
adjoint recovers one occupation integral, equal to the obstacle area;
a mixed-area identity determines the ratio from the two reciprocal
differences alone, under a weaker separation condition.

For $\gamma=0$, the common-noise converse in
Theorem~\ref{thm:stationary-lower} gives lower power $(s+1)/(s-2)$,
whereas \eqref{eq:main-new-attempts} has power $(3s/2+1)/(s-2)$.
Thus the difference of the powers is $s/(2(s-2))$; it was $2s/(s-2)$
for the retained signed-query construction. We do not identify the
stationary minimax exponent. The converse permits arbitrary-precision
adaptive short commands under a single known uniform-disk law and
charges worst-case expected attempts. The finite period assertion
continues to use the known positive nonperiod-patch margin. Exact
geometric recovery and pointwise eventual period recognition are
stated separately from that uniform finite assertion.

\subsection{Organization}
Section~\ref{sec:stationary} develops the stationary occupation inverse
and its quantitative support geometry. Section~\ref{sec:unknown-footprint}
separates two footprints when their scales are known.
Section~\ref{sec:rare-stationary} proves the faster boundary query;
Section~\ref{sec:unknown-scale} identifies an unknown ratio by integrated
collision width and, independently, by occupation area.
Section~\ref{sec:stationary-information} proves the common-noise
converse. The appendices contain the retained localized statements,
coarse reconstruction, smooth interpolation, period tests, physical
packing, finite controls and calibration converse, with their proofs.

\section{Reconstruction with a fixed, unknown stationary launch law}
\label{sec:stationary}
The calibration converse quantifies over implementation maps that may
change with the table, the command and the nominal draw. Stationary
additive noise is a different experiment. We now give a constructive
inverse for this experiment at a fixed, nonzero launch spread. The
probability density need not be known, symmetric or mean zero. Its
support, and a lower bound on its mass near the support boundary, are
calibrated apparatus data.

\subsection{A common launch law and its geometric support}
Fix a known strictly convex body $K\subset\R^2$, containing the origin
in its interior, whose support function is $C^{6,\beta}$ and whose
curvature radii have positive lower and finite upper bounds. Assume
\begin{equation}\label{eq:fixed-footprint}
                       t+\diam K<d_0.
\end{equation}
This strict inequality, the bounds on $K$, and its position in the
laboratory frame are fixed independently of the target accuracy.
Let $\gamma\ge0$ and $b_0>0$ be fixed, with the following density class
nonempty:
\begin{equation}\label{eq:jitter-class}
 \mathcal J=\left\{j\in L^1(K):\ j\ge0,\quad\int_Kj=1,\quad
 j(z)\ge b_0\dist(z,\partial K)^\gamma
                  \text{ for almost every }z\in\operatorname{int}K\right\}.
\end{equation}
For $\gamma=0$ the lower bound is $b_0$ throughout the interior. No
upper bound, derivative bound or evaluation oracle for $j$ is needed.
The law $j$ is fixed during the experiment and is independent of the
commands, their compass direction and the table. The guarantee below
is uniform over this unknown nuisance law.

At nominal center $x$, draw fresh independent $Z\sim j$ and $a\sim\rho$.
The forward preparation attempts $(x+Z,a)$; its reciprocal reverse
attempts $(x+Z+a,-a)$ with the same prescribed joint distribution and
fresh draws. Only the pooled collision bit is returned. In particular,
$Z$, the actual start, free-start status and direction-resolved means
are not observed. Durations, directions and nominal centers are exact
in this stochastic model; certified rounding of the latter is handled
below. There is no additional arbitrary implementation error.

Set
\begin{equation}\label{eq:blurred-occupation}
 v_j(x)=\int_K j(z)\one_{\mathcal O}(x+z)\dd z,\qquad
 P_C=C+(-K).
\end{equation}
The first expression is a reflected convolution when $j$ is not even.
The same reciprocal identity gives $g_j=(T-I)v_j$. Every component
$P_C$ has diameter at most $D_0+\diam K$, and distinct such components
are at distance at least $d_0-\diam K>t$. Indeed, for $c\in C,d\in D$
and $z,w\in K$, one has
$|(c-z)-(d-w)|\ge |c-d|-|z-w|$.

\begin{lemma}[Support recovery without density deconvolution]
\label{lem:noise-support}
For every $j\in\mathcal J$,
\begin{equation}\label{eq:noise-support}
 \{v_j>0\}=\bigcup_C\operatorname{int}P_C,\qquad
             p_C^{\rm lab}=p_{P_C}^{\rm lab}-p_{-K}^{\rm lab}.
\end{equation}
Consequently the exact reciprocal mean functions on a sufficiently
large fixed protected aperture determine all protected component bodies,
without knowing $j$. With exact zero tests the full period group is then
recognized as in Lemma~\ref{lem:recognition}, without a supplied positive
nonperiod margin.
\end{lemma}
\begin{proof}
The intersection $C\cap(x+K)$ has nonempty interior exactly when
$x\in\operatorname{int}(C+(-K))$. One way to see the interior assertion
is to apply the open linear map $(c,z)\mapsto c-z$ to
$\operatorname{int}C\times\operatorname{int}K$; its image is the interior
of the Minkowski difference. On that intersection $j$ is positive
almost everywhere after translation, so its integral is positive.
Outside the interior, the intersection has area zero. The separation
above makes the positive components distinct. Support functions add
under Minkowski addition, which proves the second formula.

The proof of Lemma~\ref{lem:stopped} uses only $0\le v\le1$,
separated positive components of bounded diameter, and $g=(T-I)v$.
It does not require a smooth or known convolution kernel. Applied to
$v_j$, its killed iterates recover $v_j$ uniformly on the protected
region from its exact forcing. This proves the local assertion. The
period statement applies to the recovered original bodies, not to an
assumed lattice of the noise law. These support and Minkowski facts are
classical convex geometry; see \cite{Schneider}.
\end{proof}

\subsection{A uniform lower bound near the blurred boundary}
A quantitative version of positivity is needed for a finite experiment.
For a convex body $A$, write $A_{-r}=\{x:x+r\overline B\subset A\}$.
Whenever $r$ is smaller than its uniform interior rolling radius,
$p_{A_{-r}}=p_A-r$. The identity follows by subtracting $r$ from the
support function: its curvature measure remains nonnegative, and the
resulting Minkowski summand is exactly the erosion.

\begin{lemma}[Overlap and boundary mass]
\label{lem:cap-mass}
There are $c,e_0>0$, depending only on the physical class, $K$, $b_0$
and $\gamma$, such that for $0<e<e_0$ and every protected component,
\begin{equation}\label{eq:cap-mass}
 x\in(P_C)_{-e}\quad\Longrightarrow\quad
              v_j(x)\ge c e^{\gamma+3/2}.
\end{equation}
On the complement of $\bigcup_C\operatorname{int}P_C$ the value is zero.
Both assertions are uniform over $j\in\mathcal J$.
\end{lemma}
\begin{proof}
Choose uniform interior rolling radii $r_C,r_K$ and put $r_P=r_C+r_K$.
Reduce $e_0$ so that $e_0<\min(r_C,r_K,r_P,1)/64$; after the
collar depth $d_*$ below is fixed, require also $e_0<d_*/8$.
Thus every erosion by $e/4$ used below is strictly below the relevant
rolling radius, and $K_{-e/4}$ retains a uniform positive rolling radius.
All subsequent support identities are therefore invoked only within their
valid erosion range.

First consider the unweighted overlap
$w(x)=\area(C\cap(x+K))$. Write $P=C+(-K)$.
At $x_0\in\partial P$ with outer normal $n$, the uniquely determined
support points satisfy $x_0=y-z$, where $y\in\partial C$ has normal
$n$ and $z\in\partial K$ has normal $-n$. Let $r_C,r_K$ be fixed
positive interior rolling radii, smaller if necessary than their
uniform bounds. At $x=x_0-dn$, the intersection contains the lens of
the two disks with centers $y-r_Cn$ and $y+(r_K-d)n$ and radii
$r_C,r_K$. Their centers are separated by $r_C+r_K-d$.

Here is an explicit lower bound for the lens. In tangent-normal
coordinates about $y$, choose a fixed sufficiently small $c_1>0$.
For $|u|\le c_1\sqrt d$, the top of the first disk is
$-r_C+\sqrt{r_C^2-u^2}\ge-d/4$, whereas the bottom of the second is
$r_K-d-\sqrt{r_K^2-u^2}\le-3d/4$. For small $d$ the rectangle
\[
                  |u|\le c_1\sqrt d,\qquad
                         -3d/4\le y_n\le-d/4
\]
lies in both disks. Thus $w(x)\ge c_2d^{3/2}$ throughout a fixed
inner collar of $P$. An interior point at sufficiently small distance
$d$ from $\partial P$ has the displayed form at a nearest boundary
point, so this covers the whole collar, not only prescribed normals.

The bound extends to deeper points. The function $w^{1/2}$ is concave
on $P$: the Minkowski combination of the intersections at $x_1,x_2$
is contained in the intersection at their corresponding combination,
and the planar Brunn--Minkowski inequality gives concavity. Fix a small
collar depth $d_*$. On $\partial P_{-d_*}$ the preceding bound is
$c_2d_*^{3/2}$. A chord through any point of $P_{-d_*}$ has endpoints
on this boundary. Concavity bounds $w$ there by the same lower bound.
It follows that
\begin{equation}\label{eq:unweighted-cap}
 w(x)\ge c_3\min\{\dist(x,\partial P),d_*\}^{3/2},\qquad x\in P.
\end{equation}
The constants are uniform for the small erosions of $K$ used next:
their support norms and rolling radii remain uniformly controlled.

For $x\in P_{-e}$ put $K'=K_{-e/4}$. By the support identities,
$C+(-K')=P_{-e/4}$. Hence $x$ has distance at least $3e/4$ from
the boundary of this smaller sum. Applying \eqref{eq:unweighted-cap}
to $K'$ gives
$\area(C\cap(x+K'))\ge c_4e^{3/2}$. On $K'$ the density lower
bound is $b_0(e/4)^\gamma$. Restricting the defining integral to
this intersection proves \eqref{eq:cap-mass}. For $\gamma=0$ the
same argument works, although the erosion is unnecessary.
\end{proof}


\begin{proposition}[A density-independent inverse modulus]
\label{prop:jitter-modulus}
Consider two physical experiments $(\mathcal O_0,j_0)$ and
$(\mathcal O_1,j_1)$ with $j_0,j_1\in\mathcal J$ and the same calibrated
footprint. Let $g_0,g_1$ be their reciprocal pooled-mean differences and
put $\Delta=\|g_0-g_1\|_{W,\infty}$. On a common protected patch,
for sufficiently small $\Delta$, their complete components admit a
bijection satisfying
\begin{align}
 \max_C d_H(C_0,C_1)&\le C\Delta^{1/(\gamma+3/2)},
                                      \label{eq:jitter-Hausdorff}\\
 \max_C\|p_{C_0}^{\rm lab}-p_{C_1}^{\rm lab}\|_{C^2}
    &\le C\Delta^{(s-2)/(s(\gamma+3/2))}.
                                      \label{eq:jitter-modulus}
\end{align}
The densities need not agree. On $\mathcal B_\eta$, the primitive
orbit count and rational period relations in matched generating lists
also agree below the known
locking thresholds, and matched real period bases and free areas
have error bounded by the right side of \eqref{eq:jitter-Hausdorff}.
\end{proposition}
\begin{proof}
The comparison by stopping at the zero set of the first occupation in
Section~\ref{sec:queries}, and then interchanging the two occupations,
gives
$\|v_{j_0}-v_{j_1}\|_{U,\infty}\le H_K\Delta$.
This uses the bounded separated positive components $P_C$ and not
equality of the two densities. Put
$e=(2H_K\Delta/c)^{1/(\gamma+3/2)}$, with $c$ from
Lemma~\ref{lem:cap-mass}. For a point of $P_{C_0}$ choose a point of
$(P_{C_0})_{-e}$ at distance at most $e$; the erosion identity supplies
one. Its first occupation is at least $2H_K\Delta$, hence its second
occupation is positive. The two unions of expanded components are
therefore at Hausdorff distance at most $e$, by symmetry. Protective
collars include the test points for all components under comparison.

For $e$ below a fixed fraction of the expanded-component separation,
each connected expanded component lies in the $e$-neighborhood of a
unique component of the other union. The neighborhoods are pairwise
disjoint. Applying the same argument in reverse gives a bijection of
complete components. In support coordinates their distances are at
most $e$. Subtract the same function $p_{-K}$ to obtain
\eqref{eq:jitter-Hausdorff} for the original bodies.

For completeness, on a uniformly bounded $C^{6,\beta}$ class the
interpolation estimate
\[
 \|f\|_{C^2}\le C\|f\|_\infty^{(s-2)/s},\qquad
                              \|f\|_\infty\le1,
\]
follows by convolution with an explicitly signed approximation
kernel. Starting from an even compact smooth probability kernel $\kappa$,
the linear combination of its dilates at scales $1,2,3,4$ with
coefficients $(8/5,-4/5,8/35,-1/35)$ has total mass one and vanishing
nonconstant moments through degree six. It is a signed numerical
approximation kernel, not a launch probability density. At smoothing scale $b$, its second
derivative is bounded by $C b^{-2}\|f\|_\infty$, and the derivative
remainder is $C b^{s-2}$. Choosing $b=\|f\|_\infty^{1/s}$ proves
the estimate; the first and zeroth derivatives have weaker bounds.
Apply this to the differences of matched support functions to prove
\eqref{eq:jitter-modulus}. Lemma~\ref{lem:locking} gives the final
assertion, with a matched independent-pair basis rather than a
continuously selected reduced lattice basis.
\end{proof}

\subsection{High-accuracy local queries at fixed launch spread}
The required query accuracy now tends to zero; a fixed $1/8$ error
would not locate the boundary of the positive set. A refinement of the
killed-aperture estimate avoids charging a spurious factor of the
iteration depth to every probability measurement.

\begin{lemma}[Aperture Green bound for perturbed iterates]
\label{lem:aperture-green}
Choose the fixed observation rectangle $W$ with collars protecting the
complete $P_C$ containing all targets in $U$. There are constants
$H_K,H_W,L_K$ such that the exact killed iterates for $g_j$ have error
at most $2^{-\lfloor n/L_K\rfloor}$ on $U$. If the forcing at each
state has error at most $\xi$ and each update error is at most $\zeta$,
the clipped approximate iterate satisfies
\begin{equation}\label{eq:jitter-green-error}
 \|\widehat V_n-v_j\|_{U,\infty}
       \le2^{-\lfloor n/L_K\rfloor}+H_W(\xi+\zeta).
\end{equation}
At any requested target $x$, all these iterates use only the states
$(x+t\Z^2)\cap W$, whose number is bounded by a fixed prior constant,
independently of $n$ and the target accuracy.
\end{lemma}
\begin{proof}
Take $H_K=(D_0+\diam K+t)^2/t^2$ and
$L_K=\max(1,\lceil2H_K\rceil)$. The stopped positive-component
argument proves the first assertion. For the walk killed on leaving
$W$, put $\tau_W=\inf\{n\ge0:X_n\notin W\}$ and take
$H_W=(\diam W+t)^2/t^2$. The bounded-stopping martingale calculation
gives $\sup_{x\in W}\E_x\tau_W\le H_W$.

Let $E_n$ be the pointwise absolute difference of the noisy and exact
killed iterates on the lattice through $x$. Averaging, positive part
and clipping to $[0,1]$ give
$E_{n+1}\le T^WE_n+\xi+\zeta$, with both errors zero outside $W$.
Iterating from zero yields
\[
 E_n(x)\le(\xi+\zeta)
       \sum_{k=0}^{n-1}(T^W)^k\one_W(x)
   =(\xi+\zeta)\E_x\min(n,\tau_W)\le H_W(\xi+\zeta).
\]
This bound uses exit from the entire fixed rectangle; it does not assume
that the approximate iterate knows the positive set. Combining it with
the exact error proves \eqref{eq:jitter-green-error}. A bounded rectangle
meets at most a fixed number of points of any translate of $t\Z^2$.
Transitions inside that state set are exact additions by the compass
step; transitions leaving it are killed, not wrapped or renormalized.
\end{proof}

\begin{proposition}[An indeterminate-layer query under fixed jitter]
\label{prop:jitter-query}
For each target $x\in U$, accuracy scale $0<e<e_0$ and conditional
failure allowance $0<\varepsilon<1/4$, a finite reciprocal experiment
returns the correct membership label of $\bigcup_CP_C$ whenever the
target lies outside its inner boundary layer of width $e$. In that
layer the label is unrestricted. It uses at most a fixed number of
command centers and
\begin{equation}\label{eq:jitter-query-cost}
                    C e^{-(2\gamma+3)}\log(C/\varepsilon)
\end{equation}
attempted collision bits. Its launch law has the same fixed support
$K$ at every query and every accuracy. No sample of the launch position
or of an occupation indicator is returned.
\end{proposition}
\begin{proof}
Let $b_e=c e^{\gamma+3/2}$ be the lower bound in
Lemma~\ref{lem:cap-mass}, decreasing $c$ so $b_e\le1$.
Estimate $v_j(x)$ to error at most $b_e/8$. In
Lemma~\ref{lem:aperture-green}, choose
$n=O(\log(C/b_e))$ with the first error at most $b_e/32$, and
estimate each pooled mean to error $b_e/(128H_W)$. Forcing is the
difference of two such means. Reserve update error $b_e/(64H_W)$.
These allocations give error less than $b_e/8$. There are only a
prior-bounded number of means, and their sample averages have Hoeffding
tails. A union bound gives the stated cost. Fresh preparations justify
the same assertion conditional on all previous data.

Threshold at $b_e/2$. Outside the positive components the true value is
zero, and at depth at least $e$ it is at least $b_e$. The test is therefore
correct at both kinds of points. Only the layer between them is
unrestricted. Rounding a target once by $O(e)$, followed by exact
addition of the dyadic compass step, enlarges this layer by $O(e)$;
it does not change the probability accuracy requirement. The finite
state iterates may be evaluated by rational arithmetic on the hit
counts or with the stated update reserve.
\end{proof}

\subsection{Uniform finite reconstruction and the role of calibration}
\begin{theorem}[Unknown stationary jitter on a fixed calibrated support]
\label{thm:stationary-jitter}
Fix the bounds of $\mathcal B_\eta$, a footprint $K$ satisfying
\eqref{eq:fixed-footprint}, and the nonempty density class
\eqref{eq:jitter-class}. There are constants $C,c,\nu_0>0$ and a finite
adaptive controller, independent of the unknown density $j$, such that
for $0<\nu<\nu_0$ and $0<\delta<1/4$,
\begin{equation}\label{eq:jitter-uniform-quantifier}
 \inf_{\mathcal O\in\mathcal B_\eta}\ \inf_{j\in\mathcal J}
 \Prob_{\mathcal O,j}\{\text{geometric loss}\le C\nu
              \text{ and the primitive discrete data are correct}\}
                         \ge1-\delta.
\end{equation}
The discrete data and the laboratory loss have the meanings fixed in
Section~\ref{sec:setting}. The controller uses one fixed launch support
$K$, the reciprocal compass law, and only the pooled collision bits.
It satisfies
\begin{align}
 J_\nu&\le C\nu^{-1/(s-2)}\log(C/\nu),\label{eq:jitter-centers}\\
 N_\nu&\le C\nu^{-((2\gamma+3)s+1)/(s-2)}
                  \log(C/\nu)\log(C/(\nu\delta)).\label{eq:jitter-attempts}
\end{align}
Every attempt, including a solid start or a miss, is counted. The
aperture is fixed by the priors. No localization scale of the random
start shrinks with $\nu$. For densities bounded below on $K$
($\gamma=0$), the power in \eqref{eq:jitter-attempts} is
$(3s+1)/(s-2)$.

Nominal centers have dyadic precision $O(\nu^{s/(s-2)})$ and require
$O(\log(C/\nu)+\log(C/\delta))$ digits per batch, including the
repetition count. Certified enclosures of the fixed numerical priors, including $b_0$
and $\gamma$, and certified support and derivative evaluations for $K$
are assumed when a numerical output is requested.
These are description and algorithmic assertions, not bounds on
manufacturing the launch law or calibrating its support.
\end{theorem}
\begin{proof}
The expanded bodies $P_C$ are strictly convex with uniformly bounded
$C^{6,\beta}$ support functions, positive curvature bounds, bounded
diameters and positive separation. Their curvature radii are the sums
of those of $C$ and $-K$. Choose a fixed enlarged protected aperture
for them and all coarse search brackets.

A fixed-scale version of Proposition~\ref{prop:jitter-query} gives
coarse hulls and interior centers for the expanded bodies. The proof of
Lemma~\ref{lem:coarse} applies verbatim to an indeterminate-layer label:
its deep grid nodes are retained and no node outside the expanded body
is retained; allowing an additional layer only increases the fixed
error reserve. The complete-component cutoff excludes fragments and
protects every required component. The isolated radial brackets and
positive inball therefore follow with changed prior constants.

Use $m$ angular directions, $h=2\pi/m\asymp\nu^{1/(s-2)}$, and
query at layer width $e=c h^s$. The relaxed bisection invariant of
Lemma~\ref{lem:bisection} needs only correctness outside a boundary
layer, not an exactly unbiased label or a known density. It returns
radial values of $P_C$ with error $O(e)$ in $O(\log(C/e))$ queries
per direction. Lemma~\ref{lem:radial} then gives smooth convex bodies
$\widehat P_C$ whose laboratory supports have $C^2$ error $O(h^{s-2})$
and supremum error $O(h^s)$.

Subtract the known footprint in support coordinates:
\begin{equation}\label{eq:estimated-subtraction}
                 \widehat p_C= p_{\widehat P_C}^{\rm lab}-p_{-K}^{\rm lab}.
\end{equation}
The true function $p_C+p_C''$ is bounded below by $1/\kappa_+$.
For small $\nu$, the $C^2$ error ensures
$\widehat p_C+\widehat p_C''>0$. Thus \eqref{eq:estimated-subtraction}
is the support function of an actual strictly convex body, not a formal
subtraction of sets. Its $C^2$ error is $O(h^{s-2})$ and its Hausdorff
error is $O(h^s)$. Subtracting a first harmonic for Steiner centering
preserves these orders. If a smooth output rather than a $C^{6,\beta}$
output is desired, use the finite smoothing described below.

Apply Lemma~\ref{lem:locking} to these original component bodies.
Their errors eventually lie below the known patch and bounded-denominator
thresholds, so the primitive orbit count and exact rational relations
are recovered. Real basis vectors and free area remain estimates, with
error $O(h^s)$. Positive separation and curvature yield a physical
periodic output exactly as in the proof of Theorem~\ref{thm:adaptive}.
The test is on the original bodies after subtraction; it is not an
assumption that an arbitrary noisy difference of sets preserves periods.

There are at most $Q_h=C(1+h^{-1}\log(C/h))$ adaptive local queries,
including coarse acquisition. Allocate conditional failure probability
$\delta/Q_h$ to each. Each query has a prior-bounded number of command
centers, even though its killed iteration depth grows logarithmically.
Equation~\eqref{eq:jitter-query-cost} yields total cost
\[
 C h^{-1}e^{-(2\gamma+3)}\log(C/h)\log(C/(h\delta)).
\]
Substituting $e=c h^s$ and $h\asymp\nu^{1/(s-2)}$ proves both counts.
Conditional confidence and the union bound prove
\eqref{eq:jitter-uniform-quantifier}. No property of $j$ other than its
support, lower bound and stationarity was used.

For numerical realization choose a dyadic compass step satisfying the
strict gap and one dyadic target mesh comparable to $e$. The state set
is the intersection of its translated exact compass lattice with a
fixed dyadic rectangle. Rounded radial bisection has the same mesh-floor
bound as in Theorem~\ref{thm:digital}. Sample counts here grow
polynomially in $1/\nu$, but their binary descriptions still require
only logarithmically many digits. The thresholds and count bounds can
be replaced by conservative dyadic enclosures of the fixed constants;
the receiver never needs to evaluate the unknown $j$.

For a finite smooth output, evaluate \eqref{eq:estimated-subtraction}
on an angular grid of width $b\asymp c\min(e^{1/3},\nu)$, to value
accuracy $O(b^3)$. The reconstructed supports have a uniform $C^3$
bound. Quadratic local interpolation patched with the fixed smooth
partition therefore has supremum error $O(b^3)$ and $C^2$ error
$O(b)$: Taylor remainders and value errors contribute the same orders
after up to two derivatives. These are respectively $O(e)$ and
$O(\nu)$. All support evaluations are computed from the finite radial
representation and the known footprint, not obtained through additional
collision observations. The curvature margin persists after reducing
constants. This final step is a disclosed numerical overhead, not part
of the attempted-bit count.
\end{proof}

\begin{corollary}[No fixed-jitter resolution floor in this experiment]
\label{cor:no-jitter-floor}
A fixed nondegenerate footprint satisfying \eqref{eq:fixed-footprint}
permits arbitrarily accurate reconstruction from finite collision bits,
uniformly over \eqref{eq:jitter-class}. In particular, stationary
uniform disk noise of any fixed radius $R$ with $t+2R<d_0$ has no
positive $C^2$ resolution floor in this model.
\end{corollary}
\begin{proof}
Keep $K$ and $j$ fixed and let $\nu$ decrease in
Theorem~\ref{thm:stationary-jitter}. Its sample bound is finite at every
positive accuracy. For a uniform disk, $\gamma=0$ and
$b_0=1/(\pi R^2)$ are admissible. The full density was not otherwise
needed by the construction.
\end{proof}

There is no contradiction with Theorem~\ref{thm:calibration-floor}.
The latter permits different, command- and draw-dependent
implementations in its indistinguishable worlds. Here all commands use
one stationary translate of a density with a calibrated support and
uniform boundary mass. The noise radius is not a worst-case actuator
error budget. An unknown displacement of the entire footprint, an
uncontrolled change of its support with the command, arbitrary errors
in duration or angle, and unbounded noise distributions are not included
in the new theorem. Nominal positioning and support calibration remain
separate apparatus requirements. Nor does the theorem claim a sharp
sample exponent: it proves that repetition can replace narrowing the
random launch spread in this stochastic experiment, not that every
kind of calibration can be replaced by repetition.

\section{Joint recovery of the table and an unknown launch footprint}
\label{sec:unknown-footprint}
The one-scale stationary theorem calibrates the whole convex support of
the launch law. That support can instead be an unknown of the inverse
problem when the same launch mechanism is operated at two known,
nonzero scales. The extra control is a dilation about a fixed laboratory
origin, not a position measurement. Both spreads stay fixed as the
requested accuracy tends to zero. All one-scale results remain valid
without this additional control.

\subsection{The two-scale experiment}
Fix numbers $0<\lambda_1<\lambda_2$, known to the controller. Let
$\mathcal K$ be a class of unknown strictly convex footprints with
\begin{equation}\label{eq:unknown-K-prior}
 r_K\overline B\subset K\subset R_K\overline B,\qquad
 \|p_K^{\rm lab}\|_{C^{6,\beta}}\le M_K,\qquad
 0<\rho_K^-\le p_K^{\rm lab}+(p_K^{\rm lab})''\le\rho_K^+.
\end{equation}
The constants are fixed and known, but neither $p_K$ nor its Steiner
point is supplied. The scale gap $\lambda_2-\lambda_1$ is fixed and
positive. Require
\begin{equation}\label{eq:unknown-K-separation}
                      t+2\lambda_2R_K<d_0.
\end{equation}
For each $K$, an unknown stationary density $j$ has integral one on
$K$ and satisfies
\begin{equation}\label{eq:unknown-K-density}
 j(z)\ge b_0\dist(z,\partial K)^\gamma\quad\text{a.e. in }K,
               \qquad b_0>0,\quad \gamma\ge0.
\end{equation}
Take the admissible family of pairs $(K,j)$ nonempty. An attempt at
scale $\lambda_i$ and nominal center $x$ uses $x+\lambda_i Z$,
where $Z$ is a fresh draw from $j$. The reverse attempt uses
$x+\lambda_i Z+a$ and $-a$. The compass direction is pooled as
before. Only a collision bit is returned; no realization of $Z$, impact
position, free-start status or direction-resolved outcome is supplied.
The same unknown $j$ is used throughout. Its scaled density obeys
\begin{equation}\label{eq:scaled-density-bound}
 j_i(z)=\lambda_i^{-2}j(z/\lambda_i)
 \ge b_0\lambda_i^{-(\gamma+2)}
                    \dist(z,\partial(\lambda_iK))^\gamma.
\end{equation}
Nominal centers, durations, directions and the two scale factors are
exact controls, apart from the numerical rounding specified below.
There is no additional arbitrary displacement error.

Write $v_i(x)=\int_Kj(z)\one_{\mathcal O}(x+\lambda_i z)\dd z$
and $g_i=(T-I)v_i$. The reciprocal identity and the stopped comparison
recover $v_i$ from $g_i$ on protected regions. Lemma~\ref{lem:noise-support}
gives their complete positive components as
\begin{equation}\label{eq:two-scale-components}
                 P_{i,C}=C+(-\lambda_iK),\qquad i=1,2.
\end{equation}
The controller knows uniform geometric bounds for these components,
although it does not know their support functions. Their radii of
curvature are sums, their diameters are bounded by $D_0+2\lambda_2R_K$,
and their mutual separation is at least $d_0-2\lambda_2R_K>t$.
These are all the geometric bounds needed before footprint subtraction
in the proof of Theorem~\ref{thm:stationary-jitter}.

\begin{lemma}[Matching scales without component labels]
\label{lem:scale-matching}
Put $m_K=(\lambda_2-\lambda_1)r_K$ and
$d_K=d_0-2\lambda_2R_K$. The small-scale component belonging to
$C$ satisfies
\begin{equation}\label{eq:scale-nesting}
                 P_{1,C}+m_K\overline B\subset P_{2,C}.
\end{equation}
Consequently the two collections in \eqref{eq:two-scale-components}
determine their common component correspondence without obstacle labels,
asymmetry or distinct-shape assumptions.

More precisely, suppose complete convex estimates at each scale have
Hausdorff error at most $e$, under their true within-scale correspondences.
If $e<c\min(m_K,d_K)$ for a fixed sufficiently small $c$, the Steiner
point of each estimated small-scale component lies in exactly its
corresponding estimated large-scale component. This rule recovers the
same correspondence from the estimates.
\end{lemma}
\begin{proof}
Convexity and $0\in K$ give
$\lambda_2K=\lambda_1K+(\lambda_2-\lambda_1)K$.
The latter summand contains $m_K\overline B$, proving
\eqref{eq:scale-nesting}. In particular every point of $P_{1,C}$
is at distance at least $m_K$ from the complement of $P_{2,C}$.
The distinct large-scale components are $d_K$-separated, so containment
pairs each small component with a unique large one. Every large one
contains its own small component, proving bijectivity.

The Steiner point lies in its convex body, and its change under
Hausdorff error $e$ is at most $2e$, by \eqref{eq:centering-map}.
Thus the estimated small-scale Steiner point has distance at least
$m_K-2e$ from the complement of the true large-scale component.
For compact convex bodies at Hausdorff distance $e$, support inequalities
give $A_{-e}\subset\widehat A\subset A+e\overline B$. Hence the
point belongs to the corresponding large estimate when $3e<m_K$.
It belongs to no other large estimate when $3e<d_K$. These inequalities
supply a common positive margin for certified containment tests.
No unstable choice of a nearest congruent shape is involved.
\end{proof}

All collections in this lemma consist of complete components. In a
bounded experiment take the protective collars large enough to contain
both dilations of every component needed for the period tests. The
coarse complete-component cutoff of Lemma~\ref{lem:coarse}, with the
larger diameter bound above, discards exterior fragments. Components
needed in the protected patch therefore occur at both scales; unmatched
exterior fragments are not used to infer a missing component or a period.
The enlarged aperture still depends only on the fixed priors.

\begin{theorem}[Recovery under calibrated homothety and an inverse modulus]
\label{thm:self-exact}
The exact pair of reciprocal mean functions $(g_1,g_2)$ in the
two-scale experiment determines the unknown footprint in its laboratory
frame and every protected original component. For a matched pair,
\begin{align}
 p_{-K}^{\rm lab}
   &=\frac{p_{P_{2,C}}^{\rm lab}-p_{P_{1,C}}^{\rm lab}}
                   {\lambda_2-\lambda_1},\label{eq:self-footprint}\\
 p_C^{\rm lab}
   &=\frac{\lambda_2p_{P_{1,C}}^{\rm lab}
                  -\lambda_1p_{P_{2,C}}^{\rm lab}}
                   {\lambda_2-\lambda_1}.\label{eq:self-body}
\end{align}
With exact zero tests the full period group is also determined, under
the bounded periodic presentation and without a positive patch margin.

For two admissible experiments with possibly different $K,j$ and tables,
put $\Delta=\max_i\|g_i-\widetilde g_i\|_{W,\infty}$.
For small $\Delta$, the protected components admit a correspondence with
\begin{align}
 d_H(K,\widetilde K)+\max_Cd_H(C,\widetilde C)
       &\le C\Delta^{1/(\gamma+3/2)},\label{eq:self-modulus-0}\\
 \|p_K^{\rm lab}-p_{\widetilde K}^{\rm lab}\|_{C^2}
  +\max_C\|p_C^{\rm lab}-p_{\widetilde C}^{\rm lab}\|_{C^2}
       &\le C\Delta^{(s-2)/(s(\gamma+3/2))}.
                                                   \label{eq:self-modulus-2}
\end{align}
The constants include the reciprocal of the fixed scale gap. On
$\mathcal B_\eta$, primitive discrete data lock below the retained
thresholds, and matched real period bases and free areas obey the
Hausdorff order in \eqref{eq:self-modulus-0}.
\end{theorem}
\begin{proof}
The stopped inverse first identifies the two positive-component
collections; Lemma~\ref{lem:scale-matching} identifies their pairing.
Minkowski support addition gives
$p_{P_{i,C}}^{\rm lab}=p_C^{\rm lab}+\lambda_i p_{-K}^{\rm lab}$.
Solving these two linear equations proves
\eqref{eq:self-footprint}--\eqref{eq:self-body}. The answer for $K$
is the same from every pair because the experiment uses one common
footprint. Reflection of a support argument by $\pi$ recovers $p_K$
from $p_{-K}$. This also recovers a nonzero Steiner point of $K$;
no mean-zero convention or hidden centering of the launch law was used.
Apply Lemma~\ref{lem:recognition} to the recovered original bodies.

For the modulus, the stopped comparison of Section~\ref{sec:queries}
uses only uniformly bounded, separated positive components. Thus it
compares $v_i$ and $\widetilde v_i$ even for different footprints and
different densities. Its supremum error is at most $H_*\Delta$.
The proof of Lemma~\ref{lem:cap-mass} is uniform over
\eqref{eq:unknown-K-prior}--\eqref{eq:scaled-density-bound}, so an
inner depth $e=C\Delta^{1/(\gamma+3/2)}$ has occupation greater
than $2H_*\Delta$. The erosion argument gives Hausdorff distance
$O(e)$ between the positive components at each scale, with a bijection
of complete protected components. Uniform $C^{6,\beta}$ support bounds
and the signed-kernel interpolation estimate give $C^2$ error
$O(e^{(s-2)/s})$ at each scale. This step does not subtract a supposedly
common unknown footprint.

The nesting margin makes the two within-scale correspondences consistent
with the cross-scale pairing. Indeed the small-scale Steiner point and
its counterpart remain in the corresponding large-scale bodies below
the same margin; a different pairing would put one point in two
separated large components. Now subtract the two explicit linear systems.
Their inverse has norm bounded in terms of
$(\lambda_2-\lambda_1)^{-1}$, which proves both estimates, for $K$
as well as the original bodies. Finally apply Lemma~\ref{lem:locking}
to those original bodies, after subtraction, not to the expanded union.
\end{proof}

\begin{theorem}[Finite two-scale reconstruction without a footprint oracle]
\label{thm:self-finite}
Fix the bounds of $\mathcal B_\eta$, \eqref{eq:unknown-K-prior},
the strict margin \eqref{eq:unknown-K-separation}, the two scale factors,
and $b_0,\gamma$. One finite adaptive controller, independent of $K$
and $j$, reconstructs a strictly convex footprint $\widehat K$ and a
smooth strictly convex periodic table. For all sufficiently small
$\nu>0$ and $0<\delta<1/4$, uniformly over admissible $\mathcal O,K,j$,
its geometric loss and footprint laboratory $C^2$ support error are at
most $C\nu$, and its primitive discrete data are correct, with
probability at least $1-\delta$. Its costs satisfy
\begin{align}
 J_\nu&\le C\nu^{-1/(s-2)}\log(C/\nu),\label{eq:self-centers}\\
 N_\nu&\le C\nu^{-((2\gamma+3)s+1)/(s-2)}
                \log(C/\nu)\log(C/(\nu\delta)).\label{eq:self-attempts}
\end{align}
Both random launch supports are fixed independently of $\nu$.
No footprint support or derivative evaluations and no density evaluations
are inputs. Nominal positioning is refined to
$O(\nu^{s/(s-2)})$; the exact scale factors and their common dilation
origin, the footprint envelope and regularity bounds, and the boundary-mass
constants are calibrated controls or priors. The theorem does not charge
their physical manufacture or calibration to $N_\nu$.
\end{theorem}
\begin{proof}
Run the construction of Theorem~\ref{thm:stationary-jitter} separately
at the two scales, stopping each construction before its footprint
subtraction. Every earlier step uses only occupation queries, the
component bounds and the lower boundary-mass constants; it never evaluates
$p_K$. Equation~\eqref{eq:scaled-density-bound} supplies known lower
constants at both scales. The aperture Green bound, coarse acquisition,
indeterminate-layer bisection and radial interpolation therefore give
complete convex estimates $\widehat P_{i,C}$ with support supremum
error $O(e)$ and $C^2$ error $O(h^{s-2})$, where
$h\asymp\nu^{1/(s-2)}$ and $e=c h^s$.

Use Lemma~\ref{lem:scale-matching} to pair the two finite collections.
The estimates have a common error event of probability at least
$1-\delta$ after splitting confidence between scales and queries.
On this event every pairing needed in the protected patch is correct.
For the finitely many complete pairs take the average of
\[
 q_C=\frac{p_{\widehat P_{2,C}}^{\rm lab}
                 -p_{\widehat P_{1,C}}^{\rm lab}}{\lambda_2-\lambda_1}
\]
and call it $\widehat q$. It estimates $p_{-K}^{\rm lab}$ in supremum
norm by $O(e)$ and in $C^2$ by $O(\nu)$; averaging cannot enlarge
the maximum error. Put
$\widehat p_C=p_{\widehat P_{1,C}}^{\rm lab}-\lambda_1\widehat q$.
These functions have the same two error orders for $p_C^{\rm lab}$.
Using one average makes the output footprint common to all components.
Uniform positive lower radii of curvature of $K$ and $C$ imply
$\widehat q+\widehat q''>0$ and
$\widehat p_C+\widehat p_C''>0$ for small $\nu$. Hence all these
estimates represent actual strictly convex bodies, not arbitrary formal
support differences. The known inball of $K$ persists after a fixed
reduction of its radius. The finite final smoothing of
Section~\ref{sec:stationary}, applied also to $\widehat q$, preserves
the stated orders and strict curvature.

Apply the complete-component period tests and rational locking of Lemma~\ref{lem:locking}
to the original-body estimates. All margins are met after reducing
$\nu_0$, giving a physical periodic output with the same primitive orbit
count and exact rational relations. Matched real basis vectors and free
area have error $O(e)$ and in particular $O(\nu)$. The footprint itself
is not subjected to period tests.

There are two families of radial queries. Put
$Q_h=C(1+h^{-1}\log(C/h))$. Each query uses at most
\[
                   C e^{-(2\gamma+3)}\log(CQ_h/\delta)
\]
attempts, by the fixed-aperture estimate. This proves
\eqref{eq:self-centers}--\eqref{eq:self-attempts}. Separate fresh
attempts at both scales justify conditional confidence allocation; no
coupling of the physical noise draws is assumed. Exact dyadic compass
addition, rounded radial bisection and certified arithmetic use the same
error reserves as before. The numerical constants require only the
finite priors in the theorem, not an oracle for the unknown footprint.
\end{proof}

\begin{remark}[What has, and has not, been calibrated]
\label{rem:self-contract}
This theorem replaces knowledge of an entire function $p_K$ by two
known homothetic settings and uniform geometric and mass bounds. It
does not learn the nuisance density. Exact homothety is substantive:
independently drifting unknown footprints are not covered. If the two
settings have an additional common unscaled laboratory offset $b$, the
formulas instead recover the translated table $\mathcal O-b$ and the
same $K$. Absolute table location then needs an external reference.
This follows directly from
$C-b+(-\lambda_iK)$ in \eqref{eq:two-scale-components}, and is not
a new noise floor for centered shape or for periods.
\end{remark}

\section{Boundary tests from rare pooled collisions}
\label{sec:rare-stationary}
The preceding stationary construction estimates a difference of two
collision probabilities to the accuracy of a small boundary mass.
Near an exposed boundary one can instead arrange that an exterior
test has a zero collision probability. An interior test then needs only
one observed collision. The compass direction remains unobserved.
To cover the directions at which two compass vertices maximize the same
support functional, we test a finite set of candidate vertices and take
the conjunction of their outcomes.

The compass length in this section is chosen once from the fixed
priors. Let $D_*$ be a known upper bound for the diameter of every
available launch footprint, and require
\begin{equation}\label{eq:rare-margin}
                         2t+D_*<d_0.
\end{equation}
For the calibrated one-scale experiment one may take $D_*=\diam K$;
for the two known scales of Section~\ref{sec:unknown-footprint}, take
$D_*=2\lambda_2R_K$. The extra factor of two accounts for both the
shift of a nominal center and the attempted flight. This sufficient
design condition is stronger than the one used for the earlier
reciprocal estimates. Those estimates remain valid under their stated
conditions. Whenever the footprint diameter bound is strictly below
$d_0$, a positive dyadic $t$ satisfying \eqref{eq:rare-margin} can be
chosen independently of the requested accuracy. No new output or
direction control is required.

\subsection{Coarse normal information}
The expanded components $P=C+(-K)$ have a uniform positive lower
curvature radius, a uniform support regularity bound, and a positive
separation. Fix a collar width $d_c>0$ smaller than their lower rolling
radius and such that
\begin{equation}\label{eq:rare-collar}
            d_c<t/8,\qquad D_*+2t+d_c<d_0.
\end{equation}
Reducing $d_c$ below a fixed fraction of these bounds will be harmless.
For every $y$ in this collar of a complete expanded component, there
are a boundary point $y_0$, its outward unit normal $n$, and a signed
depth $d\in[-d_c,d_c]$ such that
\begin{equation}\label{eq:rare-signed-depth}
              y=y_0-dn,\qquad n\cdot y=p_P^{\rm lab}(n)-d.
\end{equation}
Here $d>0$ inside $P$. The representation is unique in the chosen
collar: the support function $p_P-d$ still has a positive curvature
radius for every inward offset in this range, and describes the inner
parallel boundary. Exterior metric projections are unique by convexity.

\begin{lemma}[Fixed-accuracy normals for radial brackets]
\label{lem:rare-coarse-normals}
At any prescribed confidence $1-\varepsilon_c$, a prior-dependent
finite number of the earlier reciprocal queries supplies an interior
center for each protected expanded component and isolated radial
brackets containing its boundary in every radial direction. Each bracket
has a tubular neighborhood of one fixed positive prior-controlled radius
lying in the collar \eqref{eq:rare-signed-depth}. For each radial
direction the same finite data compute a rational vector $\widehat n$
satisfying
\begin{equation}\label{eq:rare-normal-error}
                         |\widehat n-n|\le1/32
\end{equation}
for the projection normal of every point in that tubular neighborhood.
The number
of attempts is at most $C\log(C/\varepsilon_c)$. For unknown
footprints the construction uses only their uniform priors.
\end{lemma}
\begin{proof}
Apply the fixed-layer version of Proposition~\ref{prop:jitter-query}
and the coarse acquisition of Lemma~\ref{lem:coarse} to the expanded
components, obtaining interior centers and isolated radial intervals.
With respect to each center, take a fixed angular mesh of width $h_c$
and refine its radial values to error $e_c\le c h_c^s$ by the same
reciprocal queries. All these parameters are fixed from the priors.
Lemma~\ref{lem:radial} yields a finite smooth radial function
$\widetilde R$ with
\[
 \|\widetilde R-R_P\|_\infty\le C h_c^s=:E_c,
 \qquad
 \|\widetilde R-R_P\|_{C^1}\le C h_c^{s-1}.
\]
Choose $h_c$ small enough that both error bounds meet the fixed
reserves below. In radial direction $u(\varphi)$, use a certified
interval containing
$[\widetilde R(\varphi)-2E_c,\widetilde R(\varphi)+2E_c]$.
Its endpoints enclose the true radial boundary, and its width can be
made arbitrarily small at this fixed stage. It remains in the original
isolated interval after reducing $h_c$ and its numerical reserve.

The normal to a positive radial graph is
\[
 \frac{R_P(\varphi)u(\varphi)
                 -R_P'(\varphi)u^\perp(\varphi)}
        {\sqrt{R_P(\varphi)^2+R_P'(\varphi)^2}}.
\]
The inball and regularity bounds make this expression uniformly
Lipschitz in $R_P,R_P'$. Substitution of $\widetilde R$ therefore
approximates the true radial boundary normal to any prescribed fixed
accuracy. The projection normals throughout the bracket have the same
property. To verify this last assertion quantitatively, let $r_*>0$
be a uniform lower rolling radius of $P$. Write $P=Q+r_*\overline B$,
using the support identity from Section~\ref{sec:stationary}. If
$x_i\in\partial P$ has outer normal $n_i$, the monotonicity of the
support-point relation for $Q$ gives
\[
 (x_1-x_2)\cdot(n_1-n_2)\ge r_*|n_1-n_2|^2,
 \qquad |n_1-n_2|\le |x_1-x_2|/r_*.
\]
The projection of any collar point lies within twice its distance
from the true radial boundary point. Hence on a tubular neighborhood
of radius $a_c$ about the bracket its normal differs from the radial
boundary normal by $O(E_c+a_c)$. Choose $h_c$ and then a fixed
$a_c>0$ so small that the whole neighborhood lies in the chosen collar
and that this error and the radial approximation error sum to less
than $1/64$. Compute a rational normal approximation with another
error less than $1/64$. This proves \eqref{eq:rare-normal-error}
uniformly on that neighborhood; its radius is independent of all
subsequent accuracy parameters.

Only a fixed number of radial values, bisections and reciprocal command
centers was used. All requested probability accuracies are fixed, so
their conditional confidence allocation costs
$C\log(C/\varepsilon_c)$ attempts. Every step precedes footprint
subtraction and depends only on the uniform bounds for the expanded
bodies. It consequently applies at either unknown-footprint scale.
\end{proof}

\subsection{A test using only positive collision records}
Put $E=\{e_1,-e_1,e_2,-e_2\}$ and, for the rational vector supplied
by Lemma~\ref{lem:rare-coarse-normals}, form the finite set
\begin{equation}\label{eq:rare-candidates}
 V(\widehat n)=
 \left\{v\in E:\ \widehat n\cdot v
       \ge\max_{w\in E}\widehat n\cdot w-\tfrac14\right\}.
\end{equation}
All these comparisons are rational. This set contains every maximizer
of $n\cdot v$ over $E$: the two normal errors cost at most $1/16$.
Moreover, every candidate satisfies
\begin{equation}\label{eq:rare-positive-normal}
 n\cdot v\ge\frac1{\sqrt2}-\frac14-\frac1{16}>\frac38.
\end{equation}
The threshold in \eqref{eq:rare-candidates} is positive, so the set
contains at most two compass vertices. In particular the construction
does not require a unique maximizing direction.

\begin{proposition}[A rare-collision boundary query]
\label{prop:rare-query}
Assume \eqref{eq:rare-margin} and condition on the successful coarse
event of Lemma~\ref{lem:rare-coarse-normals}. For a target $y$ in one
of its bracket neighborhoods, a scale $0<e<d_c$ and a conditional failure
allowance $0<\varepsilon<1/4$, pooled forward commands at at most two
shifted centers return the correct membership label whenever $y$ is
outside the expanded component or is inside it at distance at least
$e$ from its boundary. The label in the intervening inner layer is
unrestricted. The query uses at most
\begin{equation}\label{eq:rare-query-cost}
                     C e^{-(\gamma+3/2)}\log(C/\varepsilon)
\end{equation}
attempted collision bits. An exterior target is labelled zero with
probability one. No launch position, free-start indicator or compass
direction is observed.
\end{proposition}
\begin{proof}
For every $v\in V(\widehat n)$ make fresh pooled forward preparations
at $x_v=y+tv$. Mark this candidate positive if its batch contains at
least one collision. Return one precisely when every candidate is
marked positive. The random direction in each preparation still has
the prescribed uniform law on $tE$.

First exclude all other obstacles from these preparations. At the
projection in \eqref{eq:rare-signed-depth}, support addition writes
$y_0=c_0-z_0$, where $c_0\in C$ and $z_0\in K$ are the support
points with normals $n$ and $-n$. Every point of every attempted
segment has the form
\[
 c_0+(z-z_0)+tv+utw-dn,
       \qquad z\in K,\quad w\in E,\quad 0\le u\le1.
\]
Its distance from $c_0$ is at most $D_*+2t+d_c<d_0$. The physical
separation therefore excludes every component other than $C$, for
all launch offsets and all compass directions.

Suppose that $y$ is exterior, so $d<0$, and choose a true maximizing
vertex $v_*\in V(\widehat n)$. For every $z\in K$, $w\in E$ and
$0\le u\le1$, support addition gives
\begin{align*}
 n\cdot(y+tv_*+z+utw)
 &\ge p_C^{\rm lab}(n)-d+t(n\cdot v_*+u n\cdot w)\\
 &>p_C^{\rm lab}(n).
\end{align*}
The second line uses symmetry of $E$ and maximality of $v_*$, which
imply $n\cdot v_*+u n\cdot w\ge0$. Thus every segment for this
candidate misses $C$, and the preceding separation excludes all other
collisions. Its collision probability is zero. Its mark, and therefore
the conjunction, is deterministically zero. Possible boundary tangencies
are irrelevant to the asserted exterior and positive-depth labels.

Now suppose that $y$ has inner depth $d\ge e$. By
Lemma~\ref{lem:cap-mass},
\[
        \Prob\{y+Z\in C\}=v_j(y)
                    \ge c e^{\gamma+3/2}=:b_e.
\]
Decrease the fixed constant $c$ if necessary so that $b_e\le1$.
For every candidate $v$ and every $z\in K$,
\[
 n\cdot(y+tv+z)
       \ge p_C^{\rm lab}(n)-d+t n\cdot v
       >p_C^{\rm lab}(n),
\]
by \eqref{eq:rare-collar} and \eqref{eq:rare-positive-normal}.
All starts are consequently free. On the independent events
$y+Z\in C$ and $a=-tv$, the attempted segment ends in $C$ and
collides. The pooled success probability for each candidate is at
least $b_e/4$; the selected direction need never be revealed.

Use
\[
            M_e=\left\lceil\frac4{b_e}
                                     \log\frac2\varepsilon\right\rceil
\]
fresh preparations for each candidate. Its probability of having no
collision is at most $\exp(-M_eb_e/4)\le\varepsilon/2$.
A union bound over at most two candidates proves correctness at
positive depth and \eqref{eq:rare-query-cost}. Fresh preparations give
the same bound conditional on every preceding adaptive history.
No upper bound or regularity assumption on the density was used.
\end{proof}

\subsection{The improved stationary reconstruction rate}
\begin{theorem}[Stationary reconstruction by rare collisions]
\label{thm:rare-stationary}
Fix the priors of Theorem~\ref{thm:stationary-jitter} and choose a
fixed compass length satisfying $2t+\diam K<d_0$. There is a finite
adaptive controller, independent of the stationary density, whose
geometric loss is at most $C\nu$ and whose primitive discrete data
are correct with probability at least $1-\delta$, uniformly over the
same physical and density classes. For sufficiently small $\nu>0$
and $0<\delta<1/4$, its costs satisfy
\begin{align}
 J_\nu&\le C\nu^{-1/(s-2)}\log(C/\nu),
                                      \label{eq:rare-centers}\\
 N_\nu&\le C\nu^{-((\gamma+3/2)s+1)/(s-2)}
              \log(C/\nu)\log(C/(\nu\delta)).
                                      \label{eq:rare-attempts}
\end{align}
Every preparation is counted. The launch footprint and compass length
stay fixed, and all nominal centers and physical launches lie in a
prior-dependent fixed aperture. Dyadic nominal positioning to
$O(\nu^{s/(s-2)})$ suffices.

For the two known scales and unknown footprint of
Theorem~\ref{thm:self-finite}, the same conclusions and costs hold
under $2t+2\lambda_2R_K<d_0$, including laboratory $C^2$ footprint
error at most $C\nu$. No footprint-function or density evaluations
are then required. The exact scales, common dilation origin and fixed
geometric and boundary-mass priors retain their stated roles.
\end{theorem}
\begin{proof}
Perform the fixed-accuracy construction of
Lemma~\ref{lem:rare-coarse-normals}, reserving failure probability
$\delta/2$. Put $h=2\pi/m\asymp\nu^{1/(s-2)}$ and $e=c h^s$.
At each of the $m$ equally spaced radial angles, use
Proposition~\ref{prop:rare-query} for the midpoint tests of the
corresponding coarse bracket. Its rational candidate vector is fixed
for that radial direction and remains valid throughout the bracket.
The relaxed bisection invariant of Lemma~\ref{lem:bisection} applies
to the inner indeterminate layer as well. The inball converts a normal
depth of order $e$ to a radial uncertainty of order $e$. Thus
$O(\log(C/e))$ tests per direction give radial values with error
$O(e)$, without a monotonicity assumption on erroneous layer labels.

Lemma~\ref{lem:radial} gives expanded bodies $\widehat P_C$ with
support error $O(h^s)$ in supremum norm and $O(h^{s-2})$ in $C^2$.
For a known footprint, subtract $p_{-K}^{\rm lab}$ and apply the
curvature, final smoothing and period-locking arguments of
Theorem~\ref{thm:stationary-jitter}. These deterministic steps give
the stated physical output and retain the primitive conclusions.

There are at most $Q_h=C(1+h^{-1}\log(C/h))$ fine queries. Give
each one conditional failure allowance $\delta/(2Q_h)$. On the
successful coarse event, Proposition~\ref{prop:rare-query} and a
union bound give simultaneous valid labels with probability at least
$1-\delta/2$. The attempts are bounded by
\[
 C\log(C/\delta)
   +C h^{-1}e^{-(\gamma+3/2)}
                   \log(C/h)\log(C/(h\delta)).
\]
Substituting $e=c h^s$ proves \eqref{eq:rare-attempts}. At most two
shifted command centers replace each fine target; coarse reciprocal
stencils have fixed size. This proves \eqref{eq:rare-centers}.
An enlargement of the earlier fixed aperture by the fixed shifts
protects all these commands and their launch supports.

For numerical commands, round each target once to a dyadic grid of
mesh $O(e)$ and then add the exact dyadic vector $tv$. For sufficiently
small $\nu$, every rounded target lies in the fixed tubular neighborhood
of its radial bracket. The same candidate vector therefore satisfies
\eqref{eq:rare-normal-error}, even though rounding may move the target
off its radial line. Rounding changes distances to the expanded component
by $O(e)$ and hence only enlarges the indeterminate layer by that
amount. It does not turn a claimed exterior test at the rounded point
into a probability estimate. Conservative dyadic lower bounds for
$b_e$ and integer batch counts give the same costs. Their descriptions,
and those of the centers, have
$O(\log(C/\nu)+\log(C/\delta))$ digits per batch.

In the two-scale experiment apply this construction separately to
$P_{i,C}=C+(-\lambda_iK)$, using
\eqref{eq:scaled-density-bound} and the common diameter envelope
$2\lambda_2R_K$. The stronger separation assumption gives
\eqref{eq:rare-margin} at both scales. Match complete components by
Lemma~\ref{lem:scale-matching}, and use the averaged linear support
inverse and convexity verification of Theorem~\ref{thm:self-finite}.
Both scales have the same exponents and uniform prior-controlled
constants. Splitting the confidence budget and adding their costs
therefore preserves \eqref{eq:rare-centers}--\eqref{eq:rare-attempts}.
\end{proof}

For $\gamma=0$, the new upper power is $(3s/2+1)/(s-2)$.
The lower bound of Theorem~\ref{thm:stationary-lower}, specialized to
a fixed sufficiently small uniform-disk footprint and the present
compass choice, consequently gives, on a fixed physical class containing
its packing,
\begin{equation}\label{eq:rare-stationary-bracket}
 c\nu^{-(s+1)/(s-2)}
       \le N^*_{\rm stat}(\nu)
       \le C\nu^{-(3s/2+1)/(s-2)}\log^2(C/\nu).
\end{equation}
The difference between the displayed exponents is now
$s/(2(s-2))$. These bounds do not assert an exact minimax rate.
The improvement comes from a rare positive event whose absence has
an exact geometric explanation on the exterior side; it does not
follow from increasing the precision of reciprocal mean cancellation.

\section{Recovery of an uncalibrated homothety ratio}
\label{sec:unknown-scale}
Integrating a pooled forward collision mean eliminates the unknown
launch density and measures a sum of widths of the original obstacle.
Together with the two positivity supports, this scalar identifies the
ratio of two homothetic settings. An independent normalization uses
the area of a recovered occupation component and requires a weaker
separation condition.

Write $A$ for the physical footprint at the first setting, and let
$Z$ have an unknown probability density $j$ on $A$. The two settings
use the displacements $Z$ and $rZ$, where
\begin{equation}\label{eq:unk-scale-prior}
             1+g_*\le r\le R_*,\qquad g_*>0.
\end{equation}
Only these bounds, not the value of $r$, are supplied. The common
homothety origin is fixed in the nominal laboratory coordinates.
The footprint satisfies \eqref{eq:unknown-K-prior}, with $A$ in place
of $K$, and the density satisfies \eqref{eq:unknown-K-density}.
Assume $t+2R_*R_K<d_0$. Thus the two expanded-component collections
have a known positive separation greater than $t$, and their
cross-setting nesting margin is at least $g_*r_K$.

If the apparatus is originally described by unknown factors
$\lambda_1,\lambda_2$ acting on a parameter footprint $K$, then
$A=\lambda_1K$ and $r=\lambda_2/\lambda_1$. This is an operational
normalization: multiplying $K$ by a positive constant and dividing
both factors by that constant leaves the physical experiment unchanged.
The quantities recovered below are the physical first footprint and
the ratio, without assigning an additional unit to $K$.

All support functions in this section use laboratory coordinates.
For a unit vector $u$, $p_H(u)$ denotes $\sup_{x\in H}x\cdot u$,
the same support function expressed on the unit circle.

\subsection{An integrated collision width}
For a convex body $H$ write
\[
 \mathcal W(H)=p_H(e_1)+p_H(-e_1)+p_H(e_2)+p_H(-e_2),
\]
the sum of its two coordinate widths. This functional is invariant
under translation and additive under Minkowski addition. Let $A$ be
the first launch footprint, let $j$ be its probability density, and
put $P_{1,C}=C+(-A)$. Denote the first-setting pooled forward mean by
$F_1$.

\begin{lemma}[An isolated pooled flux measures width]
\label{lem:unk-scale-flux}
Suppose $2t+\diam A<d_0$. After sufficiently accurate fixed coarse
acquisition of a complete $P_{1,C}$, one can construct a bounded
domain $E$, consisting of finitely many squares of a fixed dyadic
mesh, which contains $P_{1,C}+t\overline B$ and no forward-response
support belonging to another component. Then
\begin{equation}\label{eq:unk-scale-flux}
 \Phi_C:=\int_E F_1(x)\dd x
                      =\frac t2\mathcal W(C).
\end{equation}
The domain and its area are computed from the coarse observations and
the fixed priors. Neither $j$ nor an occupation integral is an input.
\end{lemma}
\begin{proof}
For a displacement $a=tv$, $|v|=1$, set
$S_a(C)=C+[-a,0]$. A line parallel to $v$ meets $C$ in an interval,
a point, or the empty set. On every nonempty interior slice,
$S_a(C)\setminus C$ is an interval of length $t$. The transverse
projection of $C$ has length
$p_C(v^\perp)+p_C(-v^\perp)$. Cavalieri's formula therefore gives
\begin{equation}\label{eq:unk-scale-strip-area}
 \int_{\mathbb R^2}
       (\one_{S_a(C)}-\one_C)
   =t\{p_C(v^\perp)+p_C(-v^\perp)\}.
\end{equation}
Endpoint and tangency conventions affect only null sets. This is also
the segment case of the classical mixed-area formula \cite{Schneider}.

Since $t<d_0$, one segment of length $t$ cannot meet two distinct
physical components, and a start in one component cannot reach another.
Consequently the collision indicator decomposes pointwise as
\[
 B_a^{\mathcal O}
       =\sum_C(\one_{S_a(C)}-\one_C).
\]
The sum is locally finite and its summands are nonnegative. The
contribution $F_{1,C}$ obtained by averaging a single summand against
$j$ and the compass law is supported in
\[
 \bigcup_{v\in\{\pm e_1,\pm e_2\}}
       \{P_{1,C}+[-tv,0]\}
                    \subset P_{1,C}+t\overline B.
\]
Distinct expanded components are separated by at least
$d_0-\diam A$: offsets from the same footprint differ by at most its
diameter. Their displayed enlarged envelopes are therefore separated
by at least $d_0-\diam A-2t>0$. A sufficiently accurate coarse hull,
a fixed enlargement, and a sufficiently small fixed dyadic square mesh
give a domain $E$ containing the complete envelope of $C$ while
remaining disjoint from every other envelope. All error reserves use
fixed fractions of this positive separation. Thus $F_1=F_{1,C}$ on
$E$, and $F_{1,C}=0$ outside $E$.

Tonelli's theorem, translation invariance of area, and $\int_Aj=1$
now show that $\int_EF_1$ is the average of the four strip areas in
\eqref{eq:unk-scale-strip-area}. Each coordinate width occurs twice,
which proves \eqref{eq:unk-scale-flux}. No density differentiability,
upper bound, or symmetry is used.
\end{proof}

\begin{theorem}[Calibration by an integrated collision mean]
\label{thm:unk-scale-flux}
Suppose $2t+2R_*R_K<d_0$. The exact observed functions
$(g_1,g_2,F_1)$ determine the unknown ratio, the first physical
footprint in laboratory coordinates, and every protected original
component. In particular the two observed reciprocal mean pairs
suffice. Put $Q=-A$. For a matched pair $P_1=C+Q$, $P_2=C+rQ$, put
$p_D=p_{P_2}-p_{P_1}$ and use the flux $\Phi_C$ of
Lemma~\ref{lem:unk-scale-flux}. Then
\begin{equation}\label{eq:unk-scale-flux-inverse}
 a=\frac{\mathcal W(P_1)-2\Phi_C/t}{\mathcal W(D)},
 \qquad r=1+a^{-1},\qquad
 p_Q=a p_D,\qquad p_C=p_{P_1}-a p_D.
\end{equation}
The inverse is uniformly stable: expanded-support errors $e$ in
supremum norm and $u$ in $C^2$, together with flux error $m$, give
ratio and Hausdorff errors $O(e+m)$ and laboratory $C^2$ errors
$O(u+m)$. These assertions hold for sufficiently small $u+m$,
with constants depending only on the fixed priors.
\end{theorem}
\begin{proof}
The stopped inverse recovers the complete positivity supports from
$g_1,g_2$, and their fixed nesting margin gives their pairing.
Since $D=(r-1)Q$, one has $P_1=C+aD$ with $a=(r-1)^{-1}$.
Width additivity and \eqref{eq:unk-scale-flux} yield the first formula
in \eqref{eq:unk-scale-flux-inverse}; the remaining formulas follow
by support addition. The denominator is bounded below by $4g_*r_K$.
Also $a$ belongs to the fixed interval
$[1/(R_*-1),1/g_*]$. Consequently errors $e$ in the support values
and $m$ in the flux change $a$ and $r$ by at most $C(e+m)$.
The reconstructed supports then have supremum error $O(e+m)$ and
$C^2$ error $O(u+m)$. Their uniform positive curvature margins
persist at sufficiently small errors, so they represent physical
strictly convex bodies. Reflection recovers $A$ from $Q=-A$.
No first harmonic is discarded, and the original components retain
their laboratory locations.
\end{proof}

\begin{proposition}[A finite integrated collision measurement]
\label{prop:unk-scale-flux-measure}
Assume $2t+2R_*R_K<d_0$. After a prior-dependent coarse acquisition
of one complete first-setting component $P_1$, its integrated
forward mean $\Phi_C$ can be estimated to error $m$ with conditional
confidence $1-\varepsilon$ using
\begin{equation}\label{eq:unk-scale-flux-cost}
                       C m^{-2}\log(C/\varepsilon)
\end{equation}
attempted pooled collision bits. The same order bounds the number
of nominal centers. A dyadic nominal mesh of width $c m$ suffices,
uniformly over the possibly unbounded launch density.
\end{proposition}
\begin{proof}
Let $d_*>2t$ be the lower separation of the first-setting expanded
components. Choose a fixed $b<(d_*-2t)/8$. A sufficiently accurate
coarse hull determines a domain $E$, which is a finite union of
fixed dyadic squares, such that
\[
 P_1+t\overline B\subset E
       \subset P_1+(t+3b)\overline B.
\]
Choose all squares meeting an enlargement of the coarse hull;
reducing its fixed error and the fixed square size gives these
inclusions with positive margins. No other component's forward
response support meets $E$, since each such support is contained
in $P'_1+t\overline B$ and $d_*-2t-3b>0$.
Consequently $F_1=F_{1,C}$ on $E$ and
$\int_E F_1=\Phi_C$.

For $X$ uniform on $E$, one forward bit $Y$ at $X$ gives the bounded
variable $|E|Y$, whose expectation is $\Phi_C$. Independent fresh
repetitions and Hoeffding's inequality give the asserted count.
This uses the pooled forward bit itself and does not require a
direction-resolved mean.

To give the commands finite descriptions, subdivide the fixed
squares by a dyadic mesh $\zeta$ and sample their fine midpoints
uniformly. The expectation changes by $O(\zeta)$ without any
pointwise regularity assumption on $j$. Indeed, condition first
on the launch displacement and a compass direction. The hit
indicator is the difference between a convex segment sweep and
the original convex body. Intersect both with the fixed squares
of $E$. Their total boundary length and their number are uniformly
bounded. The midpoint sum differs from the integral only on fine
squares meeting these boundaries, whose total area is
$O(\zeta)$ by the convex boundary-tube estimate. The bound is
uniform in the launch displacement; averaging over its probability
law and over the four directions preserves it. Take
$\zeta\le c m$ and reserve half the error for this bias. The
confidence assertion is conditional on the coarse acquisition
and uses only subsequent fresh preparations.
\end{proof}

\subsection{Uniform finite reconstruction}
\begin{theorem}[Finite recovery without scale values]
\label{thm:unk-scale-finite}
Fix the numerical bounds of $\mathcal B_\eta$, the footprint and
boundary-mass bounds and \eqref{eq:unk-scale-prior}. Choose the fixed
dyadic compass step to satisfy $2t+2R_*R_K<d_0$. A finite adaptive
controller, supplied with the
two setting labels but neither their ratio nor a footprint or density
oracle, recovers the ratio, the physical first footprint, and a
smooth strictly convex periodic table. For all sufficiently small
$\nu>0$ and $0<\delta<1/4$, its ratio error, footprint laboratory
$C^2$ support error, and geometric loss are at most $C\nu$, and its
primitive discrete data are correct, with probability at least
$1-\delta$ uniformly over the admissible experiments.

Put
\[
                    q_\gamma=
                    \frac{(\gamma+3/2)s+1}{s-2}.
\]
The attempted-bit and nominal-center counts satisfy
\begin{align}
 N_\nu&\le
 C\nu^{-q_\gamma}\log(C/\nu)\log(C/(\nu\delta))
       +C\nu^{-2}\log(C/\delta),\label{eq:unk-scale-attempts}\\
 J_\nu&\le
 C\nu^{-1/(s-2)}\log(C/\nu)
       +C\nu^{-2}\log(C/\delta).\label{eq:unk-scale-centers}
\end{align}
Since $6<s\le7$ and $\gamma\ge0$, one has $q_\gamma>2$, so the
second term of \eqref{eq:unk-scale-attempts} is absorbed by the first.
Both physical spreads stay fixed as $\nu$ tends to zero. The finest
nominal mesh has width $c\nu^{s/(s-2)}$. The list of nominal centers,
setting labels and repetition counts has binary description length
at most
\begin{equation}\label{eq:unk-scale-control-description}
 C\left(\nu^{-1/(s-2)}\log(C/\nu)
              +\nu^{-2}\log(C/\delta)\right)
                         \log(C/(\nu\delta)).
\end{equation}
This is a bound on the digital control description. The fixed
homothety mechanism, its common origin, the known prior bounds,
and the physical realization of the nominal mesh remain apparatus
assumptions.
\end{theorem}
\begin{proof}
Run the rare-event boundary construction of
Proposition~\ref{prop:rare-query} and
Theorem~\ref{thm:rare-stationary} at both settings, stopping before
footprint subtraction. No step uses the scale value: the known
bounds on $r$ give the requisite component, curvature, separation,
and boundary-mass constants. With
$h\asymp\nu^{1/(s-2)}$ and $e=c h^s$, the two collections have
support errors $O(e)$ in supremum norm and $O(\nu)$ in $C^2$.
The fixed nesting margin pairs all complete components needed in
the protected patch. Split the confidence allowance between these
constructions and an independent integrated collision measurement.

Use one complete first-setting component and
Proposition~\ref{prop:unk-scale-flux-measure} with $m=c\nu$ to
estimate $\Phi_C$. Evaluate the four support values giving each
width sum, and apply \eqref{eq:unk-scale-flux-inverse}. The fixed
positive denominator and the stability statement of
Theorem~\ref{thm:unk-scale-flux} give $O(\nu)$ error in $r$ and
in the physical footprint. Subtract this one footprint from
every first-setting expanded support. The resulting supports have
$C^2$ and Hausdorff error $O(\nu)$, and retain strictly positive
curvature. The finite smoothing and period-locking constructions
apply after reducing $\nu_0$. Thus the output is a physical periodic
table with the required primitive data, real period coordinates,
and free area.

The two geometric reconstructions use
$C h^{-1}\log(C/h)$ local queries, each costing at most
$C e^{-(\gamma+3/2)}\log(C/(h\delta))$ attempts. The integrated
measurement adds $C\nu^{-2}\log(C/\delta)$ attempts and at most
that many centers. This proves both counts. Fresh preparations
give the required conditional estimates and a union bound gives
the stated confidence.

Choose the common fine dyadic mesh comparable to $e$ and dividing
the fixed coarse mesh in the integrated collision construction. Since $e\le c\nu$,
it also meets that construction's bias reserve. Every nominal point
lies in a fixed aperture and therefore needs $O(\log(C/\nu))$
digits. Every repetition count is polynomial in $1/\nu$ times a
logarithmic confidence factor and needs at most
$O(\log(C/(\nu\delta)))$ digits. Encoding the finite command list
therefore gives \eqref{eq:unk-scale-control-description}.
\end{proof}

\subsection{An independent area normalization}
For a component $C$, put $Q=-A$ and
\[
 P_1=C+Q,\qquad P_2=C+rQ,
 \qquad v_{1,C}(x)=\int_Aj(z)\one_C(x+z)\dd z.
\]
The subscript $C$ denotes a single complete positive component. The
stopped inverse recovers this restriction of the occupation without
knowing $j$. Fubini's theorem gives the additional invariant
\begin{equation}\label{eq:unk-scale-mass}
 M_C:=\int_{P_1}v_{1,C}(x)\dd x
     =\int_Aj(z)|C-z|\dd z=|C|.
\end{equation}
In particular, the area on the right is measured by the collision
experiment; it is not an input concerning the obstacle.

For planar convex bodies $E,F$, let $V(E,F)$ be their mixed area,
with $V(E,E)=|E|$. In support coordinates,
\begin{equation}\label{eq:unk-scale-mixed-area}
 |E|=\frac12\int_0^{2\pi}(p_E^2-(p_E')^2)\dd\theta,
 \qquad
 V(E,F)=\frac12\int_0^{2\pi}(p_Ep_F-p_E'p_F')\dd\theta.
\end{equation}
These expressions are unchanged by translation. Support addition and
the quadratic area formula are the usual mixed-area identities
\cite{Schneider}.

\begin{theorem}[An unknown scale ratio is identifiable]
\label{thm:unk-scale-exact}
Under the initial condition $t+2R_*R_K<d_0$, the exact reciprocal
differences $(g_1,g_2)$ alone determine $r$, the physical footprint
$A$ in laboratory coordinates,
and every protected original component. The full period group is then
determined by the retained exact recognition theorem.

More explicitly, match a complete pair $P_1,P_2$ by nesting and let
$D$ be the convex body with support
\begin{equation}\label{eq:unk-scale-difference}
                         p_D=p_{P_2}-p_{P_1}.
\end{equation}
Set
\[
 a_0=|P_1|,\qquad b_0^*=V(P_1,D),\qquad c_0=|D|,
 \qquad \mathcal D=(b_0^*)^2-c_0(a_0-M_C).
\]
Then $c_0>0$, $\mathcal D>0$, and
\begin{align}
 a&=\frac{b_0^*-\sqrt{\mathcal D}}{c_0}
    =\frac{a_0-M_C}{b_0^*+\sqrt{\mathcal D}},
 &r&=1+a^{-1},\label{eq:unk-scale-root}\\
 p_Q&=a p_D,
 &p_C&=p_{P_1}-a p_D.\label{eq:unk-scale-inverse}
\end{align}
All support functions in these formulas are in the laboratory frame.
No centering convention for $A$ or for its density is imposed.
\end{theorem}
\begin{proof}
Lemma~\ref{lem:scale-matching} uses only the lower nesting margin and
the separation, so it supplies the pairing without the value of $r$.
For the true pair, $D=(r-1)Q$ and $a=(r-1)^{-1}$. In particular $D$
is a full-dimensional strictly convex body containing a fixed inball.
The equation $P_1=C+aD$ and the mixed-area identity give
\[
           M_C=a_0-2a b_0^*+a^2c_0.
\]
Moreover,
\begin{equation}\label{eq:unk-scale-root-margin}
 b_0^*-ac_0=V(C,D)>0,
 \qquad \mathcal D=V(C,D)^2.
\end{equation}
Thus the true $a$ is precisely the smaller root in
\eqref{eq:unk-scale-root}. Any other positive value $a'$ for which
$p_{P_1}-a'p_D$ is the support of a convex body of area $M_C$ would
also satisfy $b_0^*-a'c_0>0$, and hence would be the same smaller
root. This proves uniqueness among physical candidates.
Formula~\eqref{eq:unk-scale-inverse} now recovers $C$ and $Q$;
reflection recovers $A$. The answer for $r$ and $A$ is common to all
matched pairs. The recognition theorem applies to the recovered
original components.
\end{proof}

\begin{proposition}[Stability of the area-normalized inverse]
\label{prop:unk-scale-stability}
Under the fixed priors above, suppose the matched expanded supports
at both settings are known with supremum error $e$ and $C^2$ error $u$, and
$M_C$ is known with error $m$. For sufficiently small $u+m$, the
smaller-root construction is well defined and gives
\begin{align}
 |\widehat r-r|+d_H(\widehat A,A)
                 +\max_Cd_H(\widehat C,C)&\le C(e+m),
                                      \label{eq:unk-scale-stability-0}\\
 \|p_{\widehat A}-p_A\|_{C^2}
       +\max_C\|p_{\widehat C}-p_C\|_{C^2}&\le C(u+m).
                                      \label{eq:unk-scale-stability-2}
\end{align}
The constants depend on the stated bounds, including $g_*$, but not
on the unknown ratio, footprint, or density.
\end{proposition}
\begin{proof}
The true $D$ has uniformly positive curvature radius. Hence the
estimated support difference remains a strictly convex support for
small $u$. On bounded convex classes area and mixed area are Lipschitz
in Hausdorff distance, so their estimated coefficients have error
$O(e)$. For example, area follows from the two parallel-body
inclusions, and mixed area follows by polarization using $|E+F|$.

The inball of $D$ and the lower perimeter bound for $C$ imply
$V(C,D)\ge c>0$. Equation~\eqref{eq:unk-scale-root-margin} therefore
keeps the discriminant uniformly away from zero. The two equivalent
root formulas have bounded derivatives throughout a fixed
neighborhood of the admissible data. Also
$1/(R_*-1)\le a\le1/g_*$. Consequently
$|\widehat a-a|+|\widehat r-r|\le C(e+m)$.
Multiplication by $p_D$, whose $C^2$ norm is uniformly bounded in the
protected frame, proves both estimates. The positive curvature
margins of the original bodies and $A$ persist, so the recovered
supports represent actual convex bodies. A single pair may be used
for $a$ and $A$, after which the same footprint is subtracted from all
other first-setting components.
\end{proof}

\subsection{A finite measurement of one component area}
The area normalization can be measured without a fine deterministic
grid of occupation queries. The following construction uses a fixed
coarse domain and a bounded weight computable before the area
measurements begin.

\begin{proposition}[Integrated occupation from pooled bits]
\label{prop:unk-scale-mass}
After a prior-dependent coarse acquisition of one complete expanded
component $P_1$, its mass $M_C$ can be estimated to error $m$ with
conditional confidence $1-\varepsilon$ using
\begin{equation}\label{eq:unk-scale-mass-cost}
                       C m^{-2}\log(C/\varepsilon)
\end{equation}
attempted pooled collision bits at the first setting. The same order
bounds the number of nominal centers. A dyadic nominal mesh of width
$c m$ suffices. All centers lie in one fixed protected aperture, and
the assertion is uniform over the possibly unbounded density $j$.
\end{proposition}
\begin{proof}
Let $d_*>t$ be the known lower separation of the expanded components.
Choose a fixed positive $b<(d_*-t)/8$. From a sufficiently accurate
coarse hull construct a domain $E$, which is a finite union of squares
of a fixed dyadic mesh $d$, such that
\begin{equation}\label{eq:unk-scale-domain}
 P_1\subset E,\qquad
 \dist(P_1,E^c)>b,\qquad E\subset P_1+3b\overline B.
\end{equation}
The mesh may be chosen to divide the dyadic compass step $t$.
Taking all mesh squares meeting a fixed small enlargement of the
coarse hull gives these inclusions, after reducing its fixed error
and $d$. In particular no other positive component meets
$E+t\overline B$. Boundaries of squares may be assigned consistently
to one neighboring square; they have measure zero.

On $E$ the first-setting forcing therefore satisfies
\[
             g_1=(T^E-I)v_{1,C}.
\]
Let $w_E=(I-T^E)^{-1}\one_E$. The killed-walk exit bound gives
$0\le w_E\le H_E$, where $H_E$ is a fixed prior constant. Because
$d$ divides $t$, the state graph of the killed walk is identical for
all starting points in the interior of any one mesh square. Thus
$w_E$ is constant on each of the finitely many squares. Its values
are the solution of a finite rational linear system with coefficients
$0,1,1/4$, and are available without any occupation observations.
The inverse exists since the walk exits the bounded domain almost
surely and has uniformly bounded mean exit time.

The compass kernel is symmetric. Integrating the killed operator on
$E$ and interchanging its finitely many translations gives
\begin{equation}\label{eq:unk-scale-duality}
 \int_E w_E g_1
   =\int_E v_{1,C}(T^E-I)w_E
   =-\int_Ev_{1,C}=-M_C.
\end{equation}
If $X$ is uniform on $E$, take one fresh forward bit $Y^+$ and one
fresh reciprocal reverse bit $Y^-$ at $X$. The variable
\[
            Z_E=-|E|w_E(X)(Y^+-Y^-)
\]
has expectation $M_C$ and absolute value at most $|E|H_E$. Independent
repetitions and Hoeffding's inequality prove
\eqref{eq:unk-scale-mass-cost} for continuous nominal sampling.
This calculation does not assume that an occupation bit is observed.

For a finite command description choose a dyadic mesh $\zeta$ dividing
$d$ and sample $X_\zeta$ uniformly from the fine square midpoints in
$E$. This changes the expectation by $O(\zeta)$, uniformly over $j$.
Here is a direct bound which does not require an upper bound or a
modulus of continuity for the density. Before averaging in the
launch displacement, each forward or reverse bit is a finite
difference of indicators of convex obstacles and their segment
sweeps, translated by the displacement. Intersect these sets with
the fixed squares of $E$, on which $w_E$ is constant. Their total
boundary length and number are uniformly bounded. A midpoint sum
and its integral agree on every fine square which does not meet
one of these boundaries. The total area of the remaining squares is
$O(\zeta)$, by the boundary-tube estimate for convex sets. The bound
is uniform in the launch displacement, so integration against the
probability density preserves it. The same argument applies to the
four pooled directions and the reciprocal shifts. Taking
$\zeta\le c m$ reserves at most half the permitted error for this
quadrature bias. Hoeffding's inequality supplies the other half.
The procedure and its confidence bound hold conditional on the
successful coarse acquisition, using fresh independent preparations.
\end{proof}

The normalization argument uses the first setting's density only
through its integral and the quantitative boundary lower bound.
It therefore remains valid if the two settings have different
stationary densities on $A$ and $rA$, provided their normalized
boundary-mass bounds are uniform. Their supports must still be
exactly homothetic. Finally, an unknown common unscaled displacement
$b$ in the commands replaces $C$ by $C-b$ throughout, leaving $r$
and the physical footprint unchanged. Translating the table and
$b$ by the same vector leaves every collision bit unchanged. The
absolute location in that extended experiment requires a laboratory
reference.

\section{An information cost of a fixed common launch density}
\label{sec:stationary-information}
The universal one-bit converse does not measure the information lost
through a fixed random launch spread. We now prove a stronger lower
bound in the stationary experiment itself. Both competing tables use
exactly the same known density and exact implementation. There is no
table-dependent noise map. The proof combines a uniform geometric
bound on the variation of every controlled collision probability with
a conditional binary-channel bound. It applies to expected stopping
times and to continuously selected adaptive commands.

\subsection{Uniform contraction of short-command probabilities}
For two periodically indexed tables with the same lattice and component
correspondences, write
$\epsilon=\max_i d_H(C_i,\widetilde C_i)$ in fixed laboratory frames.
Only the physical packing from Lemma~\ref{lem:packing} will be needed;
there is no matching problem in this definition. Let a fresh launch
law have density $j$ of compact support in a fixed ball and
$\|j\|_\infty\le J$. The density may be given in full to the controller.
Let $t_{\max}$ be a fixed upper bound on command length. Constants below
may depend on $J,t_{\max}$, the footprint envelope and the physical
priors, but not on the command center or the number of observations.

\begin{lemma}[Lipschitz contraction of collision means]
\label{lem:mean-contraction}
For the matched tables above and $0\le\epsilon\le1$,
\begin{equation}\label{eq:mean-contraction}
 \sup_{x,\,|a|\le t_{\max}}
 \left|\int j(z)B_a^{\mathcal O}(x+z)\dd z
       -\int j(z)B_a^{\widetilde{\mathcal O}}(x+z)\dd z\right|
                      \le C J\epsilon.
\end{equation}
The same bound holds for translated reverse commands and for any
parameter-independent mixture of these directions. It holds uniformly
for a family of available launch densities having a common compact
support envelope and a common $L^\infty$ bound.
\end{lemma}
\begin{proof}
For convex compact bodies $A,A'$ at Hausdorff distance at most
$\epsilon$, the inclusions
$A\subset A'+\epsilon\overline B$ and $A'\subset A+\epsilon\overline B$,
together with the planar parallel-area formula, give
\begin{equation}\label{eq:symmetric-difference-area}
 \area(A\mathbin\triangle A')
 \le\{\operatorname{Per}(A)+\operatorname{Per}(A')\}\epsilon
                                                   +2\pi\epsilon^2.
\end{equation}
The bodies and their sweeps $S_a(C)=C+[-a,0]$ have uniformly bounded
perimeters; for example each is at most $\pi$ times its diameter.
Minkowski addition by the same segment does not increase Hausdorff
distance. Thus \eqref{eq:symmetric-difference-area} bounds both solid
and swept symmetric differences by $C\epsilon$ for each matched pair.
These are standard parallel-body facts \cite{Schneider}.

Under the closed-solid convention,
$B_a=\one_{\cup_C S_a(C)}-\one_{\cup_C C}$.
The absolute difference of two union indicators is at most the sum
of absolute differences of the matched member indicators. A fixed
translate of the launch support can meet only a uniformly bounded
number of relevant solids or sweeps: their parent components meet a
ball of prior-bounded radius, and points selected from distinct parent
components are $d_0$-separated. This also bounds components whose matched
partners, rather than themselves, meet that ball. Integrate the indicator
bound and multiply by $J$ to prove \eqref{eq:mean-contraction}.
The estimate is uniform in $x$, so it also applies after the reverse
translation $x\mapsto x+a$. Integrating the bound over a commanded
mixture, or choosing among densities satisfying the same bounds,
preserves it. The noise realization is never part of the observed bit.
\end{proof}

\subsection{A binary channel with a narrow range of means}
All entropies and mutual informations in this section are in bits.
Let $h_2$ denote binary entropy, with its endpoint values defined by
continuity.

\begin{lemma}[Range bound for one binary observation]
\label{lem:binary-range}
Let $V$ have any distribution on a finite set. Conditional on $V=v$,
let $Y$ be Bernoulli with success probability $p_v$. Then
\begin{equation}\label{eq:binary-range}
                     I(V;Y)\le \max_v p_v-\min_v p_v.
\end{equation}
This holds without a positive lower bound on $p_v$ or $1-p_v$ and
also holds conditional on any prior observation history.
\end{lemma}
\begin{proof}
Put $a=\min_vp_v$ and $b=\max_vp_v$. Introduce $U$ uniform on $[0,1]$
and independent of $V$, and generate $Y=\one_{\{U\le p_V\}}$.
Revealing this auxiliary variable can only increase information. Hence
\[
 I(V;Y)\le I(V;Y,U)=I(V;Y\mid U)\le H(Y\mid U).
\]
For $U<a$ or $U>b$ the bit is deterministic independently of $V$.
For $a\le U\le b$ its conditional entropy is at most one.
Integration gives $H(Y\mid U)\le b-a$. A conditional distribution
of $V$ gives the identical proof after every history. This is an
elementary erasure comparison, not a new source-coding principle.
\end{proof}

\begin{proposition}[Stopped experiments with contracted responses]
\label{prop:stopped-contraction}
Consider $M$ parameters and an adaptive binary protocol with independent
controller seed $U_0$ and almost surely finite stopping time $\mathsf T$.
Stopping and commands are measurable from $U_0$ and preceding bits.
Suppose, after each nonterminal history and for every permitted command,
the possible response means lie in an interval of length at most $d$.
Under the uniform parameter distribution, the stopped transcript $Z$
satisfies
\begin{equation}\label{eq:contracted-transcript}
                         I(V;Z,U_0)\le d\,\E\mathsf T.
\end{equation}
If the parameters have disjoint acceptable output sets and the average
error probability is at most $\delta<1/2$, then
\begin{equation}\label{eq:contracted-Fano}
 d\,\frac1M\sum_{v=1}^M\E_v\mathsf T
                 \ge (1-\delta)\log_2M-h_2(\delta).
\end{equation}
An infinite expected cost satisfies the latter conclusion automatically.
\end{proposition}
\begin{proof}
Condition on the independent seed. Pad the transcript with a fixed
symbol after stopping and apply the chain rule up to time $n$.
A stopped history contributes zero, since stopping is determined by the
earlier history and seed. At an active history Lemma~\ref{lem:binary-range}
bounds the conditional information by $d$. Thus the information in the
first $n$ padded entries and the seed is at most
$d\sum_{k=1}^n\Prob(\mathsf T\ge k)=d\E\min(n,\mathsf T)$.
The increasing sigma fields generated by the finite padded records and
the seed generate the stopped record. Since $V$ has a finite alphabet,
conditional probabilities converge by the bounded martingale convergence
theorem; continuity and boundedness of their entropy give convergence
of the conditional entropies. Monotone convergence on the cost side
proves \eqref{eq:contracted-transcript}. Decoding by the unique acceptable
output set and the error-indicator entropy argument used in
Section~\ref{sec:sequential} give
$H(V\mid Z,U_0)\le h_2(\delta)+\delta\log_2M$.
This proves \eqref{eq:contracted-Fano}. Input coordinates, directions,
batch sizes and the stopping time supply no additional information:
conditional on the seed they are functions of the same observed record.
The launch randomness in Lemma~\ref{lem:binary-range} is an auxiliary
comparison variable, not a new physical observation.
\end{proof}

\begin{theorem}[A stationary fixed-spread lower bound]
\label{thm:stationary-lower}
For each $s=6+\beta$, there is a fixed nondegenerate subclass
$\mathcal P\subset\mathcal B_\eta$ and a single fixed, known uniform
disk launch density with the following property. Every adaptive protocol
using fresh draws of that same law, exact nominal centers, and collision
bits from commands of length at most a fixed $t_{\max}$, and having
centered component $C^2$ error at most $\nu$ with probability at least
$1-\delta$ uniformly on $\mathcal P$, satisfies, for fixed
$0<\delta<1/4$ and all small $\nu$,
\begin{equation}\label{eq:stationary-lower}
       \sup_{\mathcal O\in\mathcal P}\E_{\mathcal O}\mathsf T
                            \ge c\nu^{-(s+1)/(s-2)}.
\end{equation}
The lattice, primitive orbit count, component species and the
parameter-independent laboratory frame may all be supplied. The unknown
actual starts are not observed. Arbitrary input precision, adaptive
nominal centers and direction-specific short commands are permitted,
so the lower bound applies in particular to the pooled reciprocal
compass experiment.

The same conclusion holds if the controller may choose a launch scale
$\lambda$ in any fixed interval $0<\lambda_-\le\lambda\le\lambda_+<\infty$;
at each chosen scale the density is the known dilation of the same fixed
uniform disk law. It therefore also bounds the two-scale experiment even
after the footprint and density are disclosed in full.
\end{theorem}
\begin{proof}
Use the fixed lattice, fixed second obstacle, and support-bump family
of Lemmas~\ref{lem:packing} and \ref{lem:centered-packing}.
For every small $h$ it has $M=2^m$ parameters, $m\ge c_0/h$,
pairwise centered $C^2$ separation at least $c_1h^{s-2}$, and
\begin{equation}\label{eq:packing-small-Hausdorff}
                  \max_{\omega,\omega'}d_H(C_\omega,C_{\omega'})
                                         \le C_0h^s.
\end{equation}
The last assertion follows from the disjoint bump supports in
\eqref{eq:packing}; it is a bound in the fixed laboratory frame, not
a disclosure of a parameter-dependent translation. All other component
copies either agree or are translates of these matched bodies.

Choose a fixed $0<t_{\max}<d_0$. Choose once a disk of positive radius
small enough that its largest allowed dilation, together with $t_{\max}$,
obeys the strict separation margin. Its uniform density and all its
allowed dilations have one finite common $L^\infty$ bound and a bounded
support envelope. They also meet the stationary boundary-mass condition
with $\gamma=0$ and a fixed known $b_0>0$. The same law is used under
every parameter; indeed it is disclosed to the controller.

Lemma~\ref{lem:mean-contraction} and
\eqref{eq:packing-small-Hausdorff} imply that the range of every
permitted response probability across the packing has length at most
$d_h=C_2h^s$. This estimate is uniform over all nominal locations,
short directions, scale choices and adaptive histories. Conditioning
on a posterior supported on fewer parameters only shortens that range.
Take $h$ proportional to $\nu^{1/(s-2)}$, with its fixed factor
large enough that $c_1h^{s-2}>2\nu$. The accuracy sets are disjoint,
and Proposition~\ref{prop:stopped-contraction} gives
\[
 C_2h^s\frac1M\sum_\omega\E_\omega\mathsf T
       \ge (1-\delta)m-h_2(\delta)\ge c_3h^{-1}.
\]
The maximum expected cost is at least the average, proving
\eqref{eq:stationary-lower}. If a deterministic cap is imposed, it is
at least the same maximum expected cost. No adversarial implementation,
change of nuisance law between parameters or uncertainty in the period
structure was used.
\end{proof}

For a fixed class containing this packing, let $N^*_{\rm stat}(\nu)$
be the infimum of the worst-case expected attempts at fixed confidence
in the known uniform-disk, fixed-scale reciprocal experiment. The lower
bound applies even to a more informative direction-controlled design;
the upper bound of Theorem~\ref{thm:rare-stationary} uses only the pooled
reciprocal design, with the fixed step chosen to satisfy its strict
separation margin. Thus
\begin{equation}\label{eq:stationary-bracket}
 c\nu^{-(s+1)/(s-2)}
       \le N^*_{\rm stat}(\nu)
       \le C\nu^{-(3s/2+1)/(s-2)}\log^2(C/\nu).
\end{equation}
This is not a matching minimax theorem. It establishes a genuine
polynomial information penalty for a fixed common random spread,
beyond the universal one-bit packing bound, with a remaining gap of
$s/(2(s-2))$ between the displayed exponents. The older upper bound
of Theorem~\ref{thm:stationary-jitter}, with exponent $(3s+1)/(s-2)$,
remains valid. The new local query reduces the difference of the
upper and lower exponents by a factor of four. Neither the old factor
$2\gamma+3$ nor the new factor $\gamma+3/2$ in a sufficient boundary
query has been proved necessary for the full collision experiment. The lower
bound holds at exact nominal positions; allowing unlimited coordinate
precision cannot remove it. Conversely its constants are for a fixed
positive spread, and the proof cannot be reused with a spread tending
to zero as though those constants were uniform. This distinction
separates the stationary lower bound from both the localized upper
bound and the worst-case calibration nonidentifiability theorem.

\section{Information categories and relation to earlier inverses}
\label{sec:comparison}
The stationary results combine a geometric query reduction with an
integral normalization. The conjunction in
Proposition~\ref{prop:rare-query} produces a one-sided boundary test
from raw pooled collision bits, including at ties between compass
directions. The integrated collision strip identifies an obstacle
width sum and hence the ratio of two unknown homothetic footprints.
The area-normalized inverse supplies a second identification argument.
These reductions use the actual collision sensor and the loss specified
in Section~\ref{sec:setting}; the statistical and convex-geometric
principles used after the reductions have established antecedents.

\subsection{Boundary estimation and informative rare responses}
Active boundary estimation is an essential comparison for the finite
theorem, not merely an analogy with passive billiard spectra.
Castro and Nowak~\cite{CN} study adaptive binary classification under
boundary-fragment smoothness and noise assumptions, with matching
orders for their excess-risk criterion. Their input selects a feature
point whose response is a class label. This directly available label
is not our collision bit: a solid start and a free miss have the same
output here. Neither their observation model nor their loss identifies
our finite periodic table. We use an independent physical packing and
binary-tree proof rather than importing a classification lower bound.

Locatelli, Carpentier and Kpotufe~\cite{LCK} give an adaptive strategy
for smooth decision boundaries without prior knowledge of smoothness
and noise parameters. Their membership-query setting and local
line-search/interpolation construction make the closest algorithmic
comparison. In contrast, our smoothness is known, our error is a $C^2$
geometry norm, and the boundary label must first be derived from
pooled collision data. We claim neither a new general
bisection principle nor an advance on adaptation to unknown smoothness.
The reductions in Propositions~\ref{prop:query} and
\ref{prop:rare-query} permit this active-estimation strategy in the
physical collision experiment, at their respective costs.

Yan, Chaudhuri and Javidi~\cite{YCJ} use informative abstention
responses in active learning. Under their assumptions the rarity of
an informative response can yield an inverse-probability cost. Their
observation includes an abstention symbol. Our sensor has only its
original collision bit. The additional argument here is that translated
pooled commands and a conjunction over outward compass candidates
realize the requisite one-sided event for each boundary query. Small
occupation alone would not justify such an improvement: the individual
means in a signed reciprocal difference need not be small. The exterior
zero and interior lower bound are therefore proved for the commands
that are actually sampled.

\subsection{Localization and control precision}

The earlier scalar inverse reconstructed a uniformly accurate occupation
on a full two-dimensional grid. Its constructive upper bound was
$C\nu^{-3}\log(C/(\nu\delta))$ attempts on the same bounded
$C^{6,\beta}$ class, using $\sigma\asymp\nu^{3/2}$ and compensated
support smoothing. That theorem remains in the exact repository archive of A2 v31. The present
procedure uses $O(h^{-1})$ angular lines, a fixed-depth dependency
stencil for each queried point, and the full $C^{6,\beta}$ regularity.
It has a different adaptive command design, not a sharper analysis of
the old uniform design. We have not proved a matching lower bound
for that old design or for uniform spatial preparations.

The resource improvement is accompanied by the explicit localization
and calibration cost in \eqref{eq:precision}. Exact data remain a pair
of functionals on controlled spatial inputs. The lower bound concerns
only their finite binary realizations; it does not assert a low-information
passive invariant. The angle/time/position tolerances shrink with
localization and are not learned from the collision outcomes. Apparatus
motion, precision in representing input locations, and arithmetic
cost are distinct resources. Goldberg and Kwek~\cite{GK} treat
query-coordinate precision as an explicit resource in learning convex
polytopes by exact membership queries. The digital bound in
Theorem~\ref{thm:unk-scale-finite} follows the same accounting principle
for a different class and observation: it includes the additional
integral-measurement centers as well as the fine angular boundary
queries. It does not convert a bit-description bound into a physical
metrology or travel bound.

For completeness the exact v31 manuscript is preserved unchanged
inside the repository archive of the reviewed v33 source. It retains the arbitrary centered-law and general
nontrivial bounded-law inverses, the different-zero-set comparison,
the full fixed-aperture proof, rational interval enclosures, the
uniform-grid construction, and all earlier moment, marked-law and
count-fiber results in their original nested volumes. The more general
randomized law is an exact functional theorem; the finite local query
proved here uses the four-atom compass law. No continuous transition
law is discretized without an error budget. Rational intervals remain
conditional on forcing and survival bounds, not certificates of the
apparatus or physical prior.

The finite global conclusion also retains the known positive patch
margin and the bounded periodic presentation. Exact rational period
relations are distinct from approximate Euclidean generators. Removing
the known margin gives only pointwise eventual recognition, as described
in Section~\ref{sec:periods}. The result does not settle rigidity from
uniformly prepared count germs, exact analytic count germs, marked
length spectra or passive trajectories. These boundaries delimit the
information theorem without changing its topic or weakening any retained
mathematical assertion.

\subsection{Stationary noise versus adversarial implementation}
The new fixed-footprint theorem belongs also to support estimation with
contaminated data. Brunel, Klusowski and Yang~\cite{BKY} observe
continuous random vectors drawn from an unknown convex body and then
contaminated by additive Gaussian noise. Their geometric estimator
avoids Fourier inversion and is analyzed in Hausdorff loss. Here no
random vector or direct inside/outside response is observed: the data
are attempted collision bits at controlled centers. The fixed compact
noise support, its boundary-mass lower bound and the reciprocal reduction
are essential to our polynomial upper bound. The Gaussian observation
and its tail-based analysis are not covered by this theorem, and its
rates are not being compared as rates in a common experiment.

Support addition, inner rolling disks and the Brunn--Minkowski inequality
are classical \cite{Schneider}. Their role here is to remove an unknown
stationary density without learning that density. The positivity set of
the recovered blur, not its detailed values, identifies the footprint
sum. The retained signed high-accuracy query uses the Green bound for exit from a fixed
rectangle; its constants can be large and its finite-depth work grows
logarithmically. The area normalization uses a finite adjoint of the
same killed operator. The direct collision-strip normalization instead
uses Cavalieri's principle and width additivity; only one integrated
pooled forward mean is required. We do not claim these convex-geometric
or random-walk
principles as new abstract theorems.

The two upper bounds cannot be ranked by their attempted-bit exponents
alone. The localized model attains the smaller sample power but requires
shrinking spread and worst-case tolerance. The stationary model pays a
larger, not claimed sharp sample power while keeping its random spread
fixed and allowing the density to be unknown. The footprint and the
relative homothety factor are reconstructed in the strongest theorem.
Exact homothety about the supplied common origin, quantitative
boundary-mass and geometric bounds, per-setting stationarity and
nominal positioning remain apparatus assumptions. No theorem prices their manufacture,
extends the converse to every common noise law, or infers a passive
spectral invariant from these active experiments.

\appendix
\section{Localized reconstruction and its resource bounds}
\label{sec:localized-statements}
We retain the localized experiment, its binary converse and the earlier
stationary bounds. Their proofs follow in the remaining appendices and
in Sections~\ref{sec:stationary}--\ref{sec:stationary-information}.
The stronger stationary bounds proved in Section~\ref{sec:rare-stationary}
use the additional prior-fixed compass design stated there.

\begin{theorem}[Adaptive scalar reconstruction]\label{thm:adaptive}
Fix the numerical bounds of $\mathcal B_\eta$ and the kernel above.
There are $C,c,\nu_0>0$ such that, for $0<\nu<\nu_0$ and
$0<\delta<1/4$, a finite adaptive experiment with the reciprocal
compass commands reconstructs a smooth strictly convex periodic
presentation with geometric loss at most $C\nu$, with probability at
least $1-\delta$. It recovers the true primitive orbit count and exact
rational relations of a true period-generating list in a true independent-pair basis. The Euclidean coordinates of these vectors are estimates.

After a prior-dependent coarse acquisition, the finest localization
scale and sufficient coupled calibration are
\begin{equation}\label{eq:precision}
 \sigma=c\nu^{s/(s-2)},\qquad \ell+\tau+t\alpha\le c\sigma.
\end{equation}
The number of nominal command centers, counted with repetitions if
necessary, and the number of attempted preparations satisfy
\begin{align}
 J_\nu&\le C\nu^{-1/(s-2)}\log(C/\nu),\label{eq:centers}\\
 N_\nu&\le C\nu^{-1/(s-2)}\log(C/\nu)
                               \log(C/(\nu\delta)).\label{eq:attempts}
\end{align}
Every solid start and free miss is charged in $N_\nu$. All commands lie
in one fixed prior-controlled aperture. The reconstruction uses finite
local computations, radial bisections and finite polynomial data, not
a search over an infinite-dimensional class of physical candidates.
Its constants include the exit-time and patch-margin bounds.
\end{theorem}

Here and below a finite smooth partition of unity, with its derivatives,
is part of the fixed numerical representation. Certified evaluations
to the reserved tolerances suffice. The theorem counts preparations and
centers, not apparatus travel or bit operations for elementary functions.
Known $s$ is used; adaptivity refers to the locations of commands, not
adaptation to unknown smoothness.

\begin{theorem}[A necessary number of binary observations]
\label{thm:lower}
For each fixed $0<\beta\le1$ there is a fixed choice of the bounds
above and a nondegenerate subclass $\mathcal P\subset\mathcal B_\eta$
with the following property. Any adaptive randomized procedure whose
only table-dependent outputs are at most $N$ bits, and which recovers
the component boundaries with $C^2$ error at most $\nu$ with probability
at least $3/4$ uniformly on $\mathcal P$, must satisfy
\begin{equation}\label{eq:lower}
                         N\ge c\nu^{-1/(s-2)}.
\end{equation}
This holds even when the lattice, its primitive orbit count, component
correspondences and laboratory pose are given, and implementation is
exact. Input choices may have arbitrary precision and depend on all
previous bits. Independent controller randomness is allowed.
\end{theorem}

The construction of $\mathcal P$ is given in Section~\ref{sec:lower}.
The assertion is not a lower bound for every subclass, in particular
not for a singleton or a finite parametric family. On any bounded class containing this construction, $\eqref{eq:attempts}$ and
$\eqref{eq:lower}$ identify the power of the attempted-bit complexity.
A logarithmic gap remains. Neither result prices input precision or
converts an active experiment into a passive billiard invariant.

\subsection{Finite controls, random stopping and physical resolution}
The next results add two resource converses to the preceding constructive
inverse. The geometric regularity remains $s=6+\beta$, and all uniform
finite global decisions still require the bounded periodic class and
known positive patch margin. No new geometric response is supplied.

\begin{theorem}[Attempted bits and calibration resolution]
\label{thm:resources}
The adaptive reciprocal experiment has a realization with dyadic centers
and scales, using $O(\log(C/\nu)+\log(C/\delta))$ binary digits per
batch command and the same attempted-bit bound
\eqref{eq:attempts}. On a fixed nondegenerate physical subclass, any
procedure whose only table-dependent data are collision bits and whose
stopping decision uses only its past bits and independent randomness
has worst-case expected attempt count at least
$c\nu^{-1/(s-2)}$ at uniform confidence $3/4$.

Moreover, under table-, command- and draw-dependent bounded-position-error
implementations, uniform centered $C^2$ accuracy $\nu$ with confidence greater
than $1/2$ for short commands requires
\[
r\le C\nu^{s/(s-2)}.
\]
Here $r$ is the allowed position error, durations and directions may be
exact, and command lengths have any fixed bound less than the component
separation. Above a constant multiple of this scale,
two distinct physical tables admit identical complete adaptive bit
experiments under two possibly different admissible implementations.
No identical common-noise device is asserted. More attempted
bits do not remove this obstruction. The sufficient combined tolerance
$\ell+\tau+t\alpha\le c\nu^{s/(s-2)}$ remains as before.
\end{theorem}

Theorem~\ref{thm:digital} proves the finite-description claim and gives
a separate conservative numerical cost. Theorem~\ref{thm:expected}
proves the expected-stopping claim by a prefix-free transcript argument,
with the centering gauge explicit. Theorem~\ref{thm:calibration-floor}
proves the physical-resolution converse through a pointwise common-response
coupling of nearby convex swept bodies. Neither converse prices
apparatus production, and neither claims sharp logarithmic factors.
The deterministic-cap lower bound is retained independently: it is
not used as an unproved substitute for the sequential result.

\subsection{A contrasting experiment with fixed stationary jitter}
The preceding resolution converse does not apply to every apparatus
whose start has a bounded displacement. The next theorem recovers the
same table from a fixed, nonzero random launch spread. A known footprint
$K$ has positive curvature, a bounded $C^{6,\beta}$ support function and
$t+\diam K<d_0$. At every nominal center the apparatus draws a fresh
translate of one unknown density $j$, supported on $K$, with
$j(z)\ge b_0\dist(z,\partial K)^\gamma$ almost everywhere in its
interior. Here $b_0>0$ and $\gamma\ge0$ are known. The density is
stationary across all commands; it need not be symmetric or mean zero.
The controller does not observe the draw or its realized start.

\begin{theorem}[Fixed launch spread and unknown density]
\label{thm:fixed-spread-main}
On $\mathcal B_\eta$, under these footprint and density priors, one
finite reciprocal-compass controller, independent of $j$, reconstructs
the geometry to $C^2$ accuracy $C\nu$ and recovers the primitive discrete
data with probability at least $1-\delta$. It uses at most
\[
 C\nu^{-((2\gamma+3)s+1)/(s-2)}
               \log(C/\nu)\log(C/(\nu\delta))
\]
attempted bits in a fixed aperture. All solid starts and misses are
charged. The launch footprint is independent of $\nu$; only the nominal
command precision and probability-estimation accuracy are refined.
For a density bounded below on a fixed disk, $\gamma=0$ and the power
is $(3s+1)/(s-2)$. No optimality of this sample exponent is asserted.
\end{theorem}

Theorem~\ref{thm:stationary-jitter} proves the full statement.
The mechanism is geometric rather than Fourier deconvolution. The
reciprocal means recover a blurred occupation whose positive components
are $C+(-K)$. Subtraction of the known support function of $-K$ recovers
$C$. Lemma~\ref{lem:cap-mass} controls the small probabilities near
this expanded boundary. Proposition~\ref{prop:jitter-modulus} also
gives a H\"older inverse modulus uniformly over different unknown
stationary densities. Theorem~\ref{thm:calibration-floor}, by contrast,
uses two possibly different table-, command- and draw-dependent bounded
implementation maps. These are different uncertainty classes for the
same binary collision sensor, not conflicting universal resolution laws.

\subsection{Unknown footprints and the cost of a common random spread}
The fixed-spread theorem still supplies an entire footprint support
function. We next remove that input in a two-scale experiment, and give
a noise-specific converse without adversarial implementation maps.

\begin{theorem}[Calibrated homothety and stationary information cost]
\label{thm:self-main}
\textup{(i)} Let two fixed positive scales $\lambda_1<\lambda_2$ of one
unknown convex launch footprint be available, with known dilation origin,
uniform inball, envelope and $C^{6,\beta}$ curvature bounds, and
$t+2\lambda_2R_K<d_0$. Suppose one unknown stationary density on the
footprint satisfies the known lower bound
$j(z)\ge b_0\dist(z,\partial K)^\gamma$. On $\mathcal B_\eta$,
finite pooled reciprocal collision bits recover both the table and the
footprint in laboratory $C^2$ error $C\nu$, and the primitive discrete
data, at confidence $1-\delta$. The attempted-bit upper bound is
\eqref{eq:self-attempts}, with the same exponent as the calibrated
one-scale bound. No footprint-function or density oracle is supplied.
Both spreads remain fixed; nominal positioning and exact homothetic
scale control are separate resources.

\textup{(ii)} For one fixed known uniform-disk launch law and exact
implementation, a nondegenerate fixed physical class requires
\[
                  \sup_{\mathcal O}\E_{\mathcal O}\mathsf T
                         \ge c\nu^{-(s+1)/(s-2)}
\]
at fixed confidence. This holds even with known lattice and component
species, arbitrary-precision adaptive nominal positions, and
individually chosen short directions. It also holds when any launch
scale in a fixed interval bounded away from zero can be selected.
The noise law is the same under all competing tables, and no actual
start is observed. This lower bound does not match the stationary
upper exponent; the exact gap is stated in \eqref{eq:stationary-bracket}.
\end{theorem}

Sections~\ref{sec:unknown-footprint} and
\ref{sec:stationary-information} prove these assertions. The first
uses the nested geometric supports at the two scales to identify
component correspondence and to separate the footprint from the body
by support addition. The second proves that the means of every allowed
bit vary by only $O(h^s)$ on a physical packing with $\log_2M$ of
order $h^{-1}$. A conditional binary-channel inequality then applies
through arbitrary stopping rules. Thus the statistical cost is not
inferred from the older worst-case calibration converse or from the
number of bits alone.

\section{From reciprocal collisions to local occupation queries}
\label{sec:queries}
This section supplies the observation reduction used by every adaptive
query. The random walk is computational: its exit time is not an
observed collision record or an apparatus stopping signal.

\begin{lemma}[Reciprocal balance and calibration]\label{lem:balance}
With $\chi=\one_{\mathcal O}$, outside null boundary sets,
\begin{equation}\label{eq:balance}
 B_a(q)-B_{-a}(q+a)=\chi(q+a)-\chi(q).
\end{equation}
Consequently the difference $g_\sigma$ of the two pooled nominal means
satisfies
\begin{equation}\label{eq:forcing}
 g_\sigma=(T-I)u_\sigma,\qquad
 Tv(x)=\tfrac14\sum_{a\in\{\pm te_1,\pm te_2\}}v(x+a).
\end{equation}
For either command, coupled errors of size
$r=\ell+\tau+t\alpha\le\sigma$ change its mean by at most $Cr/\sigma$.
The constant is uniform over the protected aperture, components and
nominal directions, including tangencies.
\end{lemma}
\begin{proof}
If both endpoints are free, a segment and its reverse either both hit
or both miss. If exactly one endpoint is solid, the direction starting
there gives zero and the other direction hits. If both are solid both
bits are zero. This proves \eqref{eq:balance}. Averaging in $q,a$ and
commuting translation with convolution proves \eqref{eq:forcing}.

For a convex body put $S_a(C)=C+[-a,0]$. The hit bit is
\[
 B_a(q)=\one_{\cup_C S_a(C)}(q)-\chi(q).
\]
The displacement error is at most $\tau+t\alpha$. Support inequalities
show that membership can change only if the nominal start lies in
an $r$-tube of $\partial C$ or $\partial S_a(C)$, including the position
error. The part of a convex boundary tube in a disk of radius $\sigma$
has area at most $C(\sigma+r)r$. One proof integrates boundaries of
outer parallel bodies and inner erosions by the coarea formula; their
portions in the disk have length at most $2\pi\sigma$, by monotonicity
of the perimeter of convex sets. Degenerate levels follow by a limit.
Only a bounded number of components contribute: each has a point in
a disk of radius $\sigma+t+r$, and such points from distinct components
are $d_0$-separated. Multiplication by the density bound
$\sigma^{-2}\|\kappa\|_\infty$ proves the assertion. The tube containment
is pointwise for every allowed error, so the coupling need not have
mean zero. An exceptional coupling event adds its probability.
\end{proof}

Put $D=D_0+2\sigma_0$, $H=(D+t)^2/t^2$ and
$L_* =\max(1,\lceil2H\rceil)$. The positive set of $u_\sigma$ lies in
the $\sigma$-enlarged components. They have diameter at most $D$ and
mutual distance greater than $t$.

\begin{lemma}[Stopped inverse in one aperture]\label{lem:stopped}
For the compass walk $X_n=x+a_1+\cdots+a_n$ and
$\tau_u=\inf\{n\ge0:u_\sigma(X_n)=0\}$ one has
\begin{equation}\label{eq:survival}
 \sup_x\E_x\tau_u\le H,\qquad
 \sup_x\Prob_x(\tau_u>n)\le2^{-\lfloor n/L_*\rfloor}.
\end{equation}
Let a bounded target region $U$ contain the complete protected bodies
and their search brackets. Choose a fixed bounded $W$ containing
$U+\overline B_{D+t+1}$. On $W$ use the killed transition
$T^Wv=T(\one_Wv)$ and set iterates to zero outside $W$. For
$V_0=0$, $V_{n+1}=\max(0,T^WV_n-g_\sigma)$,
\begin{equation}\label{eq:bellman-error}
 0\le u_\sigma-V_n\le2^{-\lfloor n/L_*\rfloor}\quad\hbox{on }U.
\end{equation}
Uniform forcing error $\xi$ and numerical update error $\zeta$ add
at most $n(\xi+\zeta)$ if approximate iterates are clipped to $[0,1]$.
\end{lemma}
\begin{proof}
Until its first zero, a path stays in one enlarged component, since
one step cannot reach another. For $S=\tau_u\wedge n$,
$|X_S-x|\le D+t$. The bounded stopping identity for the martingale
$|X_n-x|^2-nt^2$ gives $t^2\E S\le(D+t)^2$. Monotone convergence
proves the mean bound. Markov's inequality at $L_*$ and the Markov
property at successive blocks prove the geometric tail.

On $W$, $T^Wu_\sigma\le Tu_\sigma$, so $0\le V_n\le u_\sigma$
and the iterates vanish on its zero set. Starting in $U$, the complete
containing component and every first-exit location lie in $W$. Along
this stopped path the killed transition has no boundary discrepancy.
For $e_n=u_\sigma-V_n$, the recursion gives $e_{n+1}\le Te_n$ before
exit, and hence
\[
 e_n(x)\le\E_x[u_\sigma(X_n)\one_{\{\tau_u>n\}}].
\]
This proves \eqref{eq:bellman-error}. Maximum, averaging and clipping
are nonexpansive in supremum norm; induction proves the error bound.
No periodic wrap or renormalization of a boundary row is used.
\end{proof}

The stopped argument also compares two physical occupations with
unrelated zero sets. Indeed, for $w=u-\widetilde u$ and
$h=(T-I)w$, bounded stopping at $\tau_u\wedge n$ gives
\[
 w(x)=\E_xw(X_{\tau_u\wedge n})
                 -\E_x\sum_{j<\tau_u\wedge n}h(X_j).
\]
At exit the terminal value is $-\widetilde u\le0$. On survival it
is at most one. Letting $n$ increase and then interchanging $u$ and
$\widetilde u$ yields $\|u-\widetilde u\|_{U,\infty}
\le H\|(T-I)(u-\widetilde u)\|_{W,\infty}$. Thus exact reciprocal
functionals at all sufficiently small scales determine the local
occupation and, by approximate-identity convergence and regular
closedness, the physical set. This retained exact conclusion requires
neither analyticity nor a supplied zero set. The martingale and killed
Green mechanisms are classical; see \cite{LawlerLimic}.

\begin{proposition}[A finite local query from pooled bits]
\label{prop:query}
For any adaptively chosen $x\in U$ and $\sigma\le\sigma_0$, an
occupation estimate with error less than $1/8$ can be computed using
only a prior-bounded number of reciprocal command centers in $W$.
At conditional confidence $1-\varepsilon$ it uses at most
$C\log(C/\varepsilon)$ attempted preparations. A sufficient calibration
is $\ell+\tau+t\alpha\le c\sigma$. Thresholding the estimate at
$1/2$ gives the correct membership label outside the $\sigma$-tube
of the component boundaries. Its label inside that tube is unrestricted.
\end{proposition}
\begin{proof}
Fix $n_*=6L_*$. To evaluate $V_{n_*}(x)$, its dependency locations
are exactly among
\begin{equation}\label{eq:diamond}
 x+t(i,j),\qquad |i|+|j|\le n_*.
\end{equation}
Intersect these with $W$ and put zero at locations outside it. Evaluate
the recursion backwards through shrinking diamonds, initializing
$V_0=0$. Induction on the layer shows equality with the full killed
iterate at the target. There are $O(n_*^2)$ centers and $O(n_*^3)$
updates, both independent of $\sigma$. No fine background spatial grid,
interpolation of the transition, or growing aperture is involved.

At each center estimate each pooled mean to error at most
$1/(1024n_*)$, and impose the same bound on its calibration bias.
The forcing error is then at most $1/(256n_*)$. Reserve update error
at most $1/(256n_*)$. Lemma~\ref{lem:stopped} gives total error at most
$1/64+1/256+1/256<1/8$. Hoeffding's inequality and a union bound over
the fixed finite number of means give the preparation bound. The
constant includes $n_*$ and may be large. The proof holds conditional
on past data because the new attempts are fresh. Inside a component
at boundary distance greater than $\sigma$ the average is one;
outside all enlarged components it is zero. Thresholding therefore
has the stated property. A membership label is a computed conclusion
from these bits, not an extra response furnished by the sensor.
\end{proof}

The full launch region may be taken to contain
$W+\overline B_{t+\sigma_0+1}$, which includes reverse positions and
kernel supports. For at most $Q$ adaptive queries, allocating conditional
failure probability $\delta/Q$ to each gives joint correctness with
probability $1-\delta$. It suffices to know a deterministic upper bound
on $Q$; repeated centers may be resampled. No independence of the
adaptively selected locations themselves is asserted.

\section{Adaptive recovery of complete component boundaries}
\label{sec:boundary}
We first acquire a coarse, fixed-resolution description. Its purpose
is to find interior centers and isolated radial brackets, not to supply
the final accuracy. This avoids both a supplied free reference point
and a bisection that might accidentally encounter a different obstacle.

\begin{lemma}[Coarse hulls and interior centers]\label{lem:coarse}
A prior-dependent finite number of queries from
Proposition~\ref{prop:query}, at one fixed scale $\sigma_c>0$, recovers
all complete components needed in the protected patch to Hausdorff
error $e_c\le3\sigma_c$. It supplies, for each such component $C$, a
polygon $H$ and a computable center $o$ for which
\begin{equation}\label{eq:inner-ball}
 \overline B_r(o)\subset C\cap H\subset\overline B_{R_1}(o)
\end{equation}
with fixed positive $r,R_1$. For every direction $v$, it supplies a
radial interval of fixed length at most $C e_c$ containing the boundary
point of $C$ on $o+\R_+v$. The interval is disjoint from the
$\sigma$-enlargements of all other components for every sufficiently
small fine scale $\sigma$.
\end{lemma}
\begin{proof}
The curvature upper bound supplies an interior rolling radius
$r_* =1/\kappa_+$. In support coordinates this follows from
$p+p''\ge r_*$: subtracting the support function of the radius-$r_*$
disk leaves a support function with nonnegative curvature measure.
In particular $C=C_{-2\sigma_c}+2\sigma_c\overline B$ when
$2\sigma_c<r_*$. Choose $\sigma_c<\min(r_*/8,d_0/40)$ and a square
grid of covering radius $\rho_c\le\sigma_c/8$ on a surrounding
square with collar larger than $2D_0+2$.

At its nodes retain the estimated occupations at least $1/2$, join
retained nodes at distance at most $4\sigma_c$, and take a convex hull
for each connected cluster. Every retained node is within $\sigma_c$
of a unique body; all nodes of $C_{-\sigma_c}$ are retained. For any
retained node assigned to $C$, choose a point in $C_{-2\sigma_c}$ at
distance at most $3\sigma_c$, and then a nearest grid node. That node
is retained and joined to the first. Convexity connects any two such
deep nodes by a segment; partitioning it into lengths at most $\rho_c$
and replacing partition points by nearest nodes yields a retained
path with successive lengths at most $3\rho_c$. Different bodies
cannot join because $d_0-2\sigma_c>4\sigma_c$. The hull thus lies in
$C+\sigma_c\overline B$, whereas every point of $C$ is within
$2\sigma_c+\rho_c$ of a retained deep node. This proves the error.
Discard clusters whose hulls are within $D_0+1$ of the outer square
boundary. Fragments cannot survive this cutoff, and the chosen collar
protects every required complete component.

Compute a largest inscribed disk of $H$ by its finite half-plane
inequalities, to any fixed reserved precision. Since
$d_H(H,C)\le e_c$, the support inequalities imply
$C_{-e_c}\subset H$ and $H_{-e_c}\subset C$. A disk of radius $r_*$
in $C$ therefore gives an inscribed disk of radius at least $r_*-e_c$
in $H$. The computed center, with the radii reduced by the error and
numerical reserve, satisfies \eqref{eq:inner-ball} for a fixed $r>0$.
Diameter bounds supply $R_1$.

Write $R_H(v),R_C(v)$ for the radial functions about $o$. Since both
bodies contain $\overline B_r(o)$,
\[
 H+e_c\overline B\subset o+(1+e_c/r)(H-o),
\]
and the analogous inclusion holds for $C$. The two Hausdorff
inclusions give $|R_H(v)-R_C(v)|\le C e_c$ uniformly. The polygonal
radial value is computable from its sides. Use a slightly enlarged
bracket about it, allowing numerical error. By decreasing the fixed
$\sigma_c$, its length and every point's distance to $C$ are less
than $d_0/4$. Separation excludes every other component from the
bracket and its sufficiently small kernel neighborhood.
\end{proof}

\begin{lemma}[Bisection with an indeterminate boundary layer]
\label{lem:bisection}
On each bracket in Lemma~\ref{lem:coarse}, the derived label of
Proposition~\ref{prop:query} may be arbitrary, and may depend on the
past, only when
\begin{equation}\label{eq:radial-ambiguity}
                       |a-R_C(v)|\le A\sigma
\end{equation}
for a fixed $A$. After $k$ midpoint tests, ordinary bisection returns
an estimate with error at most $A\sigma+2^{-k-1}w_0$, where $w_0$
is the initial bracket length. In particular $O(\log(C/\sigma))$
queries recover $R_C(v)$ to $O(\sigma)$.
\end{lemma}
\begin{proof}
Translate $o$ to zero. The inner ball implies
\[
 (1-\sigma/r)C\subset C_{-\sigma},\qquad
 C+\sigma\overline B\subset(1+\sigma/r)C.
\]
Together with $R_C\le R_1$, these show that outside
\eqref{eq:radial-ambiguity} the label is the true radial inside/outside
label. Other components were excluded from the bracket. A label one
at radius $a$ implies $a\le R_C+A\sigma$; a label zero implies
$a\ge R_C-A\sigma$. If the current endpoints are $l,u$, midpoint
updates consequently preserve the relaxed invariant
$R_C\in[l-A\sigma,u+A\sigma]$. This does not require $R_C$ to remain
in $[l,u]$, or the noisy labels to be monotone. The width halves at
every step. Its midpoint has the asserted error.
\end{proof}

\begin{lemma}[Radial regularity and stable interpolation]
\label{lem:radial}
Under \eqref{eq:inner-ball} and the physical support priors, the
$2\pi$-periodic radial function $R_C(\varphi)$ has a uniformly bounded
$C^{6,\beta}$ norm. Let $h=2\pi/m$ for an integer $m\ge9$. From
radial values at the $m$ equally spaced angles, each with error at most
$e$, one can construct a smooth function $\widehat R$ such that
\begin{equation}\label{eq:radial-interpolation}
 \|\widehat R-R_C\|_{C^j}
           \le C(h^{s-j}+e h^{-j}),\qquad 0\le j\le3.
\end{equation}
For $e\le C_0h^s$ and sufficiently small $h$, the curve
$o+\widehat R(\varphi)(\cos\varphi,\sin\varphi)$ is the boundary of
a smooth strictly convex body $\widehat C$, and
\begin{equation}\label{eq:geometric-interpolation}
 d_H(\widehat C,C)\le Ch^s,\qquad
 \|p_{\widehat C}^{\rm lab}-p_C^{\rm lab}\|_{C^2}\le Ch^{s-2}.
\end{equation}
Here uncentered support functions are compared in the same laboratory
frame; centering preserves these orders.
\end{lemma}
\begin{proof}
Use the support function $p$ relative to $o$, so $p\ge r$. In normal
angle $\theta$, the boundary is $p n+p' n^\perp$, its radial angle is
\[
 \varphi=\theta+\arctan(p'/p),\qquad
 \frac{\dd\varphi}{\dd\theta}
             =\frac{p(p+p'')}{p^2+(p')^2}>c>0.
\]
The inverse $\theta(\varphi)$ is uniformly $C^{5,\beta}$. Although
the expression $R=\sqrt{p^2+(p')^2}$ initially gives only five
ordinary derivatives, differentiation in radial angle cancels the
curvature factor:
\begin{equation}\label{eq:radial-cancellation}
                   R'=R(p'/p)\circ\theta.
\end{equation}
Its right side is uniformly $C^{5,\beta}$. Thus $R$ is uniformly
$C^{6,\beta}$, without losing a derivative from the support prior.

At each angular node interpolate the seven consecutive sampled
values by a polynomial of degree at most six in a local lift of the
angle. Combine these polynomials with a fixed translated smooth
partition of unity supported on neighborhoods of size $O(h)$.
On each such neighborhood the normalized Vandermonde inverse is
bounded. A degree-six Taylor polynomial of the true function is
reproduced exactly, and its remainder and derivatives through order
three are $O(h^{s-j})$. Errors in the values contribute $O(eh^{-j})$.
Derivatives falling on the partition have size $O(h^{-j})$ and are
absorbed by the corresponding lower-order remainders. The overlap
number is fixed. This proves \eqref{eq:radial-interpolation}, including
across the angular cut by periodic lifts.

A positive radial function is strictly convex when
\[
       F=R^2+2(R')^2-RR''>0.
\]
For the true curve this expression has a uniform positive margin,
by the inradius, curvature and diameter bounds. It persists for
$\widehat R$ by the $C^2$ estimate. The two radial graphs are uniformly
$O(h^s)$-close, which proves the Hausdorff bound.

To check the stronger support norm, set
$\theta=\varphi-\arctan(R'/R)$, now viewing the normal angle as a
function of radial angle. Its derivative is
$F/(R^2+(R')^2)>c$. At corresponding angles the support value and
radius of curvature are
\[
 p=\frac{R^2}{\sqrt{R^2+(R')^2}},\qquad
 p+p''=\frac{(R^2+(R')^2)^{3/2}}{F}.
\]
These expressions, and the tangential support derivative, depend
Lipschitzly on $R,R',R''$ on the bounded set with these margins.
The inverse normal-angle maps differ by $O(\|\widehat R-R\|_{C^1})$.
Their composition changes the true expressions by the same order,
using the uniform $C^3$ bound; the reconstructed functions also have
such a bound by \eqref{eq:radial-interpolation}. Thus the support
error through two derivatives is $O(h^{s-2})$. Translation to the
laboratory frame adds the same first harmonic to both supports. The
Steiner point is a bounded linear functional of the support, so its
error is $O(h^s)$ and centering preserves the conclusions.
\end{proof}

\begin{proposition}[Adaptive patch reconstruction]\label{prop:patch}
For a fixed protected patch, complete component boundaries can be
recovered with the errors in \eqref{eq:geometric-interpolation} using
at most $C h^{-1}\log(C/h)$ local occupation queries. The finest
localization is $\sigma=c h^s$. With probability at least $1-\delta$
the attempted-bit cost is at most
\begin{equation}\label{eq:patch-cost}
           C h^{-1}\log(C/h)\log(C/(h\delta)).
\end{equation}
The number of components, the initial coarse queries and every
local Bellman stencil size are bounded solely by the fixed prior.
\end{proposition}
\begin{proof}
On the simultaneous accuracy event, obtain the hulls and centers in
Lemma~\ref{lem:coarse}. Separation and the fixed aperture bound their
number. For each body use $m$ equally spaced radial directions and
their isolated coarse brackets. Lemma~\ref{lem:bisection} uses
$O(\log(C/h))$ midpoint queries per direction to return values with
error $e\le C\sigma$. Lemma~\ref{lem:radial} gives the claimed
geometric estimates. The total number of queries is bounded by a
deterministic $Q_h=C(1+h^{-1}\log(C/h))$. Allocate failure probability
$\delta/Q_h$ to every query, including the coarse ones, before the
experiment. Conditional confidence and the union bound in
Proposition~\ref{prop:query} give the simultaneous event. Its
preparation cost is \eqref{eq:patch-cost}. Locations chosen using
coarse or previous fine data do not alter this argument.

For a finite numerical realization, polygon intersections, midpoint
locations and radial values are enclosed with error $O(\sigma)$.
The estimate \eqref{eq:radial-interpolation} absorbs those value errors.
Represent the output by its finite polynomial coefficients and the
fixed smooth partition functions. They, their derivatives and the
normal-angle maps can be evaluated to any positive prescribed error
by finite operations on compact intervals with the stated derivative
bounds. This constructs an actual smooth convex output, without
assigning a false smoothness bound to the initial polygons. No
infinite-dimensional candidate minimization is performed.
\end{proof}

\section{Finite binary descriptions of the adaptive controls}
\label{sec:digital}
The active protocol specifies a distribution by a finite command, rather
than transmitting an exact real coordinate. We separate this digital
resource from the physical production of a localized launch distribution.
The latter remains a calibrated apparatus primitive; it is not obtained
by observing the collision bits.

\begin{theorem}[A dyadic controller and its resource bounds]
\label{thm:digital}
On the same bounded class $\mathcal B_\eta$, choose the compass length
$t>0$ and the fixed coarse numerical reserves dyadic, while retaining
all strict separation margins. The experiment of
Theorem~\ref{thm:adaptive} can be implemented with dyadic nominal
centers and dyadic localization scales. Its accuracy, deterministic
attempt bound and exact discrete outputs are unchanged up to constants.
Put
\[
 b=\lceil C_0+\log_2(1/\sigma)\rceil,
 \qquad B=1+\lceil\log_2(C/\nu)\rceil
              +\lceil\log_2(C/\delta)\rceil .
\]
The finest center mesh is $2^{-b}\le c\sigma$, with
$b=O(\log(C/\nu))$. Each batch command has a binary description of
length at most $CB$, and the total transmitted command description is
at most $CJ_\nu B$. The finite boundary representation has length
at most $C\nu^{-1/(s-2)}B$. These bounds include batch repetition
counts. They do not count the analog random numbers internal to the
prescribed launch distribution as table-dependent observations.

For rational fixed numerical priors and rational $\beta$, one may
choose the fixed smooth partition and its evaluators so that a
conservative digital bit-operation bound is
\begin{equation}\label{eq:digital-work}
              C(N_\nu+J_\nu+\nu^{-1/2})B^c
\end{equation}
for an absolute finite exponent $c$. This includes boundary output,
primitive arithmetic and free-area evaluation to error $C\nu$.
It is not a sharp computational complexity bound. For general fixed
$\beta$ the message bounds hold with certified dyadic scale enclosures;
a bit-operation bound additionally depends on their computation cost.
\end{theorem}
\begin{proof}
\emph{One exact dyadic stencil.}
Take the fine localization to be a dyadic number comparable to $c h^s$,
where $h=2\pi/m\asymp\nu^{1/(s-2)}$. Choose a dyadic rectangle $W$
with all the protective collars of Lemma~\ref{lem:stopped}. Round a
requested target $x$ once to a dyadic center $x'$, with
$|x-x'|\le C2^{-b}$. Use the exact stencil
\[
                         x'+t(i,j),\qquad |i|+|j|\le n_*.
\]
For $b$ greater than the fixed denominator exponent of $t$, all stencil
points are on the same dyadic grid. Membership in $W$ is tested by
exact rational comparisons. There is no independent rounding of
transition targets, hence no error in the killed compass identity.
The reverse nominal distribution is exactly the translated joint law:
it attempts $(q+a,-a)$ for the same specified law of $(q,a)$, not an
independently rounded reverse start. Fresh attempts are still used.

The query estimates $u_\sigma(x')$. Its derived label at the intended
point $x$ is therefore correct outside the boundary tube of width
$\sigma+C2^{-b}$. Physical position, duration and angle errors receive
an independent portion of the reserve in
Proposition~\ref{prop:query}. Thus rounding and actual calibration
together preserve its fixed accuracy bound. The means are ratios of
integer hit counts to a deterministic batch size. Taking that size to
be a common integer upper bound for all queries makes every local
average and Bellman update an exact rational computation. Clipping and
comparison with $1/2$ are also exact. The depth and stencil size are
prior constants.

\emph{Rounded adaptive geometry.}
The coarse grid and its finite polygons have a prior-bounded number
of rational vertices. A dyadic interior center with a reserved fixed
inball can be found by finite search; squared distances to rational
sides test the inball inequalities. The strict inball reserve ensures
that a sufficiently fine fixed search grid succeeds. Approximate
$\cos(2\pi j/m)$ and $\sin(2\pi j/m)$ with certified error at most
$c2^{-b}$. Polygonal radial values and brackets are computed with the
same reserve. Consequently every implemented spatial midpoint is
within $C2^{-b}$ of the intended radial point.

For completeness, rounding a radial midpoint by at most $d$ does not
accumulate $kd$ after $k$ queries. The relaxed interval invariant from
Lemma~\ref{lem:bisection} holds with a boundary layer of width
$A(\sigma+d)$. Its interval widths satisfy
\[
                     w_{k+1}\le w_k/2+d,
 \qquad w_k\le2^{-k}w_0+2d.
\]
Stop before asking for resolution finer than the mesh. With
$d=C2^{-b}\le c\sigma$, $O(\log(C/\sigma))$ queries give radial
values of error $O(\sigma)$, even when labels in the enlarged layer
are chosen adaptively. Round the final values to the same mesh.
Lemma~\ref{lem:radial} absorbs these errors and still gives actual
smooth strictly convex output bodies. The known patch margin and
bounded-denominator gaps allow the same exact period decisions.

\emph{Description lengths.}
All centers lie in a fixed rectangle. Their dyadic integer numerators
therefore use $O(b)$ bits. A batch command consists of a forward/reverse
flag, two such coordinates, a scale exponent, the fixed compass/kernel
identifier, and a repetition count of order $\log(C/(\nu\delta))$.
The last count requires only $O(\log B)$ binary digits. Fixed-width
fields give the asserted $CB$ bound without a real-valued control
field. There are at most $CJ_\nu$ batches, allowing repetitions.
The receiver interprets a fixed kernel identifier as the stipulated
localized random preparation, not as a table-dependent data stream.

Store the dyadic radial values, the dyadic centers and the integers
$j,m$. In normalized coordinates the seven-point interpolation matrix
has fixed rational entries, so its coefficients have $O(B)$-bit
rational descriptions. Local angles and partitions are generated by
a fixed program from $j,m$; no table-dependent exact trigonometric
constant is stored. There are $O(m)$ coefficient blocks, proving the
boundary-description bound.

\emph{A conservative operation count.}
For the computational assertion choose an explicit smooth bump,
for example $\exp[-1/(1-x^2)]$ on $|x|<1$, zero outside, and normalize
a finite-overlap family of its translates to make the fixed partition.
The denominator has a positive fixed lower bound. To the required
precision, the bump and its first three derivatives are evaluated by
rational series on bounded arguments; near its endpoints an explicit
exponential tail bound permits certified truncation to zero. The
trigonometric functions, $\pi$, square roots and logarithms needed
here admit rational series and bisection with a polynomial number of
bit operations in $B$. For rational $s$, dyadic powers comparable to
$h^s$ can equivalently be bracketed by raising rational quantities to
fixed integer powers. Thus all scale choices are effective in this
computational subcase.

Processing bit counts, the constant-depth rational Bellman updates,
rounded bisection and the fixed-degree coefficient blocks costs
$C(N_\nu+J_\nu)B^c$. One need not optimize global numerical work to
obtain a bound: the reconstructed support functions have uniform
$C^2$ bounds. Composite trapezoidal quadrature with
$O(\nu^{-1/2})$ nodes computes their Steiner integrals, and the area
integrals $\tfrac12\int\widehat R^2$, to error $O(\nu)$. A support
value is evaluated by inverting the uniformly monotone normal-angle
map; bisection with residual enclosures uses $O(B)$ steps. Each
partition evaluation involves only a fixed number of neighboring
coefficient blocks. The uniform derivative and convexity margins
justify all these tolerances.

The number of protected components and candidate period pairs is a
prior constant. Support-defect comparisons need only a fixed precision
smaller than the known patch and rational-separation margins, whereas
the real period coordinates and areas are evaluated to $O(\nu)$.
Their mesh tests, bounded-denominator rational search and two-dimensional
Hermite arithmetic are finite with polynomial cost in $B$. Including
the preceding quadratures proves \eqref{eq:digital-work}. The larger
$\nu^{-1/2}$ term is a disclosed conservative numerical overhead; it
is not silently identified with the attempted-bit exponent.
\end{proof}

\subsection{Separate resource axes}
Write $m_\nu=\nu^{-1/(s-2)}$ and $L_\nu=\log(C/\nu)$. The following
accounting refers to the construction, not to a measured apparatus.
\begin{center}
\small
\begin{tabular}{@{}p{0.31\textwidth}p{0.64\textwidth}@{}}
\toprule
Resource & Bound or explicit status\\
\midrule
Table-dependent outputs & $N_\nu\le C m_\nu L_\nu\log(C/(\nu\delta))$ bits;
all attempts, including zeros, are charged.\\[3pt]
Command centers/batches & $J_\nu\le C m_\nu L_\nu$, counted with repetitions.\\[3pt]
Digital controls & $O(B)$ bits per batch; $O(J_\nu B)$ in total.\\[3pt]
Spatial localization & $\sigma\asymp\nu^{s/(s-2)}$; the localized kernel
is a prescribed preparation primitive.\\[3pt]
Physical calibration & $\ell+\tau+t\alpha\le c\sigma$ is sufficient.
Section~\ref{sec:calibration-floor} gives a matching necessary position-error power.\\[3pt]
Aperture & Fixed by the geometric prior, including reverse starts and
protective collars.\\[3pt]
Numerical bit operations & The conservative bound \eqref{eq:digital-work}
under its explicit effective-prior assumptions.\\[3pt]
Motion, setup, calibration production & No physical time, energy or
cost law is assumed. No bound on these resources is inferred from $N_\nu$.\\
\bottomrule
\end{tabular}
\end{center}
Digital description length is not mechanical resolution. In particular,
a finite word specifying a narrow kernel does not prove that a device
can produce that kernel cheaply. The converse calibration theorem
below quantifies a necessary tolerance in the observation model,
without inventing a law for the cost of achieving it.

\section{Period recognition from the reconstructed patch}
\label{sec:periods}
This section retains the protected-patch argument for repeated shapes.
It acts on the acquired geometry, not on an assumed list of integer
copy labels. We include the proof to specify exactly which discrete
information the adaptive experiment recovers.

Put $B=1+L_0$ and $Q_a=[-a,a]^2$. The aperture used above is chosen
large enough to recover complete components with centers in $Q_{10B}$,
together with their protective collars. For two components define
\[
 d_*(C,D)=\|z_C-z_D\|_\infty+\|p_C-p_D\|_\infty,
\]
where the centered support functions are compared in the laboratory
directions, without an independent rotation. Define the patch defect
\begin{equation}\label{eq:defect}
 \mathcal D(w)=\max_{C:z_C\in Q_B}\ \max_{\epsilon\in\{-1,1\}}
                     \inf_D d_*(C+\epsilon w,D).
\end{equation}
The infimum ranges over physical components. Distant components cannot
minimize it. The known nonperiod margin is the requirement
\begin{equation}\label{eq:patch-margin}
 \begin{gathered}
 w=z_D-z_C,\quad z_C\in Q_{3B},\ z_D\in Q_{5B},\quad w\notin\Pi
       \quad\Longrightarrow\quad\mathcal D(w)\ge\eta>0.
 \end{gathered}
\end{equation}
This is a margin for whole-patch translations, not a separation between
all individual shapes. Every fixed table in $\mathcal B$ has some
positive such margin, since there are finitely many candidate pairs.
The margin is not claimed to be uniformly inferable from finite bits.

\begin{lemma}[Exact patch recognition]\label{lem:recognition}
One has $\mathcal D(w)=0$ if and only if $w\in\Pi$. For
$\|w\|_\infty\le8B$, this zero test uses only components centered in
$Q_B$ and $Q_{9B}$. The group $\Pi$ is a lattice with
\begin{equation}\label{eq:index-bound}
 [\Pi:L\Z^2]\le r_0,\qquad\covol\Pi\ge V_*=V_0/r_0.
\end{equation}
Choose any component centered in $Q_{2B}$ as a root. Its center
differences to components centered in $Q_{4B}$ that pass the zero
test generate $\Pi$ over $\Z$.
\end{lemma}
\begin{proof}
Each orbit modulo $L\Z^2$ has a representative centered in
$L[-1/2,1/2]^2\subset Q_B$. A zero defect gives both translates of
every such representative as physical components. Translating by
$L\Z^2$ gives $\mathcal O+w\subset\mathcal O$ and
$\mathcal O-w\subset\mathcal O$; together these imply equality.
The converse is immediate. The centers of a matching pair in the
bounded test are in $Q_{9B}$.

Map a coset of $\Pi/L\Z^2$ to the component orbit of the translated
root. This is injective: a nonempty compact body is not invariant
under a nonzero translation. Thus the quotient has at most $r$ elements.
A finite extension of a full-rank lattice is a lattice, giving
\eqref{eq:index-bound}. The two columns of $L$, and representatives
of all these period cosets, can be chosen with norm less than $B$
in the maximum norm. Translating the root by these vectors gives
targets in $Q_{3B}$, strictly within the cutoff $Q_{4B}$. They are
included in the accepted list and generate $\Pi$. Conversely each
accepted vector is a true period.
\end{proof}

\begin{lemma}[Stable discrete decisions]\label{lem:locking}
Suppose the protected complete components have matched Hausdorff
error at most $e$. On $\mathcal B_\eta$, for sufficiently small $e$
depending on the full prior, finite tests recover the true primitive
orbit count and exact rational relations of a true accepted generating
list in a true independent-pair basis. They recover a matched period
basis and primitive free area to error $O(e)$.
\end{lemma}
\begin{proof}
Hausdorff error $e$ gives Steiner-center error at most $2e$ and
centered-support error at most $3e$. Use approximate source centers
in $Q_{2B}$, a root there and root targets in $Q_{4B}$. For a measured
difference $\widehat w$, its corresponding true difference has error
at most $4e$. Test both signs on all selected sources, minimizing over
targets with approximate centers in $Q_{10B}$. The center contribution
to each comparison changes by at most $8e$ and the support contribution
by at most $6e$. True periods therefore have measured defect at most
$14e$. For a nonperiod, its true endpoints lie in $Q_{3B},Q_{5B}$,
so \eqref{eq:patch-margin} supplies a witness of defect at least $\eta$.
All witnesses centered in $Q_B$ are selected. Restricting targets cannot
lower the true infimum. Its measured defect is at least $\eta-14e$.
Reserve another $6e$ for numerical enclosures and use threshold $\eta/2$,
with $20e<\eta/4$. This accepts exactly the true periods. Strict cutoff
collars ensure that every generator in Lemma~\ref{lem:recognition}
is present. The same tests on pairs of sources give their primitive
orbit relation: a period translating the first center to the second
must translate the first body to the second body, since Steiner points
lie in the interiors and different components are disjoint. All
primitive orbits are represented, so their number is recovered.

Accepted true vectors have a prior bound $U_0$ on their Euclidean
lengths. A nonzero determinant of two such periods is an integer
multiple of $\covol\Pi\ge V_*$. Determinant perturbations are
$O(U_0e)$; hence independence is decided below this gap. Select an
independent true pair with matrix $D$. Its generated sublattice has
index
\[
       n=|\det D|/\covol\Pi\le Q:=\max(1,\lceil U_0^2/V_*\rceil).
\]
Every coordinate of $D^{-1}w_j$ has reduced denominator dividing $n$
and a prior-bounded numerator. The inverse-matrix identity bounds its
estimation error by $C_*e$. Distinct rationals of denominators at most
$Q$ are separated by at least $Q^{-2}$. Finite search locks the exact
coordinates when $C_*e<1/(3Q^2)$.

Let $q$ be their common denominator, which divides $n$. A column
Hermite basis $H_0$ of the integer span of $qI_2$ and $qD^{-1}w_j$
gives
\begin{equation}\label{eq:period-basis}
                         \Pi=(DH_0/q)\Z^2.
\end{equation}
The integer group contains $q\Z^2$, so its Hermite entries have a
prior bound. Replacing $D$ by its estimate gives a matched true basis
with error $O(e)$. We have recovered exact integer relations, not
exact real coordinates. The primitive free area is
\[
             A_* =\covol\Pi-\sum_{j=1}^{r_*}\area(C_j).
\]
Uniform perimeter bounds and the parallel-body inclusions give area
error $O(e)$ for each matched component, proving the last assertion.
Support suprema and center integrals used in the tests admit certified
finite angular meshes, because the bounded bodies have a uniform
support Lipschitz bound. Their numerical errors are charged above.
\end{proof}

\begin{proof}[Proof of Theorem~\ref{thm:adaptive}]
Choose a sufficiently large fixed target square so that every component
needed in Lemma~\ref{lem:locking} and every coarse bracket is protected.
Separation bounds their number. Choose the fixed observation and launch
apertures as in Lemma~\ref{lem:stopped}. Apply
Proposition~\ref{prop:patch} with an integer angular mesh satisfying
$h\asymp\nu^{1/(s-2)}$ and a sufficiently small constant factor.
Its reconstructed bodies have matched $C^2$ support error $O(\nu)$
and Hausdorff error $e=O(h^s)$. For small $\nu$, all thresholds in
Lemma~\ref{lem:locking} hold. It gives exact primitive discrete data
and period-basis error $O(h^s)$. Repeating one recovered body per
primitive orbit under this basis gives a physical periodic presentation.
Indeed the inverse basis remains bounded, so only a bounded number of
integer translates can enter a bounded fundamental patch. The positive
separation and curvature margins then persist. The reconstructed
bodies have area error $O(h^s)$ and the covolume has the same order;
hence the primitive free-area error is also $O(h^s)$, in particular
$O(\nu)$.

The query count of Proposition~\ref{prop:patch} and its fixed local
stencil size give \eqref{eq:centers}. Equation~\eqref{eq:patch-cost}
gives \eqref{eq:attempts}, since $s$ is fixed. The calibration bias
allocation in Proposition~\ref{prop:query} gives \eqref{eq:precision}.
Every preparation has already been charged there. All steps are finite,
and the output is a finite smooth representation with certified
numerical tolerances. This proves the theorem.
\end{proof}

Without a known $\eta$, successive refinements give the corresponding
pointwise eventual decisions for each fixed table, but not a uniformly
valid observable stopping certificate. The exact functional inverse
combined with Lemma~\ref{lem:recognition} needs no supplied positive
margin: it tests zeros exactly. This distinction, as well as the
bounded periodic presentation itself, is retained. No inference of
crystallinity from an arbitrary nonperiodic configuration is made.

\section{A physical packing and a binary-transcript bound}
\label{sec:lower}
The lower bound counts the table-dependent binary outputs, not the
precision of the controller's choices. It therefore applies to more
powerful binary command designs as well as the reciprocal protocol.
Its geometric construction has fixed periods and a fixed positive
patch margin, so period uncertainty is not the cause of the bound.

\begin{lemma}[A uniform physical packing]\label{lem:packing}
For $s=6+\beta$ there is a fixed class $\mathcal P\subset\mathcal B_\eta$
and constants $c_0,c_1>0$ such that, for all sufficiently small $h$,
$\mathcal P$ contains $2^m$ tables, where $m\ge c_0/h$, whose variable
component support functions have pairwise $C^2$ distance at least
$c_1h^{s-2}$. All have the same known primitive lattice and two
primitive component orbits. Their correspondences and pose can be
supplied without reducing this separation.
\end{lemma}
\begin{proof}
Fix the lattice $20\Z^2$. Put one variable component near the unit
disk at the origin, and a fixed disk of radius two at $(6,0)$, and
repeat both by this lattice. Choose a nonzero smooth function
$\psi$ supported in $(-1/3,1/3)$ with $\psi''(0)\ne0$.
Take angular centers $\theta_j$ in a fixed interval with mutual
spacing at least $h$, with $c/h\le m\le C/h$. For
$\omega\in\{0,1\}^m$, define the uncentered support of the variable
body by
\begin{equation}\label{eq:packing}
 p_\omega(\theta)=1+a h^s
                   \sum_{j=1}^m\omega_j\psi((\theta-\theta_j)/h),
\end{equation}
periodically extended, where $a>0$ is a fixed sufficiently small
number. At most one summand is nonzero at each point. Its derivative
of order $k\le6$ is bounded by $Ca h^{s-k}$. The sixth derivative
has uniformly bounded $\beta$-H\"older seminorm. Within a support
this follows by rescaling; between supports it follows from their
separation and the vanishing at their edges. The same estimate follows
by comparison to the intervening zero point for two arbitrarily close
points on different sides of an edge. Thus the $C^{6,\beta}$ bound
is uniform over all $h,m,\omega$.

For small $a$, $p_\omega+p_\omega''$ lies in a fixed positive interval.
The bodies are smooth and strictly convex, with all separation,
diameter, lattice and curvature bounds independent of $h$. Subtracting
their Steiner first harmonic also preserves the uniform support prior.
The two component species cannot be translated into each other:
their diameters lie in disjoint fixed intervals. Every period must
therefore send a variable component to another variable component,
forcing that period to lie in $20\Z^2$. Conversely every such vector
is a period. This proves primitiveness and the stated orbit count.

Center differences of the same species are true periods. For any
cross-species difference, the translated source species cannot match
a target of the other species with small support error, because their
widths differ by a fixed amount. A target of the same species has
center discrepancy bounded below by a fixed positive number: the
cross-species offset $(6,0)$, perturbed by the variable Steiner point,
stays a positive distance from $20\Z^2$. There are representatives
of both species in $Q_B$. These two alternatives give a uniform
positive defect for every nonperiod candidate in
\eqref{eq:patch-margin}. Hence one fixed $\eta>0$ works throughout
the construction. Taking the union of these tables over all small
$h$ defines a single class $\mathcal P$ with fixed numerical bounds.

If $\omega\ne\omega'$, choose an index where they differ and evaluate
the second derivative at its center. All other summands vanish there,
so the absolute difference is $a|\psi''(0)|h^{s-2}$. This proves the
separation in the laboratory support norm. The variable component,
lattice and fixed component may be identified in advance; none of
these disclosures identifies its bump vector. No analytic smoothness
bound is used or claimed for this packing.
\end{proof}

\begin{lemma}[Binary decision trees with randomized commands]
\label{lem:bits}
Let $M$ parameters have pairwise disjoint acceptable output sets.
A procedure with at most $N$ binary table-dependent outputs, arbitrary
adaptive inputs and parameter-independent randomization has average
success probability at most $2^N/M$ under the uniform distribution on
these parameters. The binary response probabilities may be arbitrary
functions of the parameter and chosen input.
\end{lemma}
\begin{proof}
Represent the controller randomization by an independent seed. Each
binary response can be generated by comparing another independent
uniform random variable with its conditional response probability.
This constructs the joint adaptive experiment for every parameter on
one common seed space. Fix all these independent seeds. Commands are
then functions only of the preceding binary transcript, and the output
is a function of a binary tree with at most $2^N$ leaves, after padding
short transcripts if necessary. An output at a leaf is acceptable for
at most one parameter, by disjointness. Thus at most $2^N$ parameters
can succeed for these fixed seeds. Average over the uniform parameter
and then integrate over the seeds. The bound follows. The numerical
values of selected commands give no additional leaves, since they are
already determined by the seed and transcript.
\end{proof}

\begin{proof}[Proof of Theorem~\ref{thm:lower}]
In Lemma~\ref{lem:packing} take
$h\asymp\nu^{1/(s-2)}$ with its constant large enough that the support
separation exceeds $2\nu$. The $\nu$-accuracy output sets are disjoint.
Lemma~\ref{lem:bits}, with $M=2^m$, gives average success at most
$2^{N-m}$. Uniform success at least $3/4$ therefore requires
$N\ge m+\log_2(3/4)\ge c\nu^{-1/(s-2)}$ for small $\nu$.
This argument allows more input control than the upper-bound procedure
and assumes perfect calibration. It remains valid for the latter
procedure with all-attempt bookkeeping. A protocol that reveals a
free-start test or a continuous impact location has a different output
alphabet and is not covered by the assertion.
\end{proof}

The loss in this theorem is the supplied laboratory-frame loss used by
the active controller. It is already enough to establish a matching
power for Theorem~\ref{thm:adaptive}, in the same loss. We do not silently
replace it by an unlabelled quotient loss in the lower bound, nor extend
a deterministic cap $N$ to an expected stopping-time budget without
another argument. No lower bound for the logarithms or for confidence
dependence is asserted.

For fixed confidence, on a nondegenerate bounded class containing
$\mathcal P$, the two results give
\begin{equation}\label{eq:matched-power}
 c\nu^{-1/(4+\beta)}\le N^*(\nu)
          \le C\nu^{-1/(4+\beta)}\log^2(C/\nu),
\end{equation}
where $N^*$ is the least uniform deterministic cap on attempted bits,
and constants in the accuracy definition are absorbed by rescaling
$\nu$. The matching power comes from one-dimensional boundary
smoothness measured two derivatives below the prior. The physical
work needed for the upper bound is that these boundary queries can
be implemented from the stipulated collision bits without additional
geometric outputs or uncharged free preparations.

\section{The lower-bound gauge and sequential observations}
\label{sec:sequential}
The deterministic-cap result above is retained. We now make its gauge
explicit and prove a separate result for expected stopping times. No
real-valued stopping signal or unobserved free-start test is allowed.

\subsection{A parameter-independent laboratory frame}
Let $n(\theta)=(\cos\theta,\sin\theta)$. For a laboratory support $p$
write
\begin{equation}\label{eq:centering-map}
 z(p)=\frac1\pi\int_0^{2\pi}p(\theta)n(\theta)\dd\theta,
 \qquad p^0=p-z(p)\cdot n.
\end{equation}
The disclosed information in Lemma~\ref{lem:packing} consists of the
fixed laboratory axes, $20\Z^2$, the names of the two component species,
and the fixed radius-two disk at $(6,0)$. These data are identical for
all packing parameters. In particular, the unknown Steiner point of
the variable body is not disclosed. ``Pose given'' does not mean a
parameter-dependent translation or rotation of that body.

For the variable component define acceptable output sets by
\begin{equation}\label{eq:accuracy-sets}
 \mathcal U_\omega(\nu)=
 \{\widehat C:\ \|p^0_{\widehat C}-p^0_\omega\|_{C^2}\le\nu\}.
\end{equation}
This centered-shape loss alone is enough for the lower bound. The
stronger laboratory loss may also include
$|z_{\widehat C}-z(p_\omega)|$ and errors of the fixed component.
No independent rotation of each output is taken in either loss.

\begin{lemma}[The packing survives centering]\label{lem:centered-packing}
For the family \eqref{eq:packing}, distinct parameters satisfy
\[
 \|p^0_\omega-p^0_{\omega'}\|_{C^2}\ge c h^{s-2}
\]
for all sufficiently small $h$, uniformly over the packing. The
centered physical support bounds and the patch margin remain uniform.
\end{lemma}
\begin{proof}
The disjoint angular supports give
$\|p_\omega-p_{\omega'}\|_\infty\le C a h^s$, independently of
$m$. Equation~\eqref{eq:centering-map} implies
$|z(p_\omega)-z(p_{\omega'})|\le2Ca h^s$. The $C^2$ norm of the
subtracted first harmonic is therefore $O(a h^s)$. At a differing
bump center the uncentered second-derivative difference has magnitude
$a|\psi''(0)|h^{s-2}$. Since $h^s=o(h^{s-2})$, at least half this
difference survives centering. A first harmonic is bounded in every
fixed smoothness norm by its amplitude, so the centered prior is also
preserved. The separation and period arguments of
Lemma~\ref{lem:packing} already allow these $O(a h^s)$ center shifts.
Thus the disclosed data contain no shape-dependent centering information.
\end{proof}

\subsection{A stopped transcript carries at most its expected length}
A stopping rule is admissible when, conditional on parameter-independent
controller randomness, stopping, further commands and the eventual
output depend only on the preceding bits. Waiting times or a physical
success flag are not additional data. The number $\mathsf T$ counts
all attempted preparations, including zeros from solid starts.

\begin{lemma}[Sequential binary information bound]\label{lem:sequential}
Let $M\ge2$ parameters have disjoint acceptable output sets. If an
admissible randomized procedure succeeds with probability at least
$1-\delta$ for each parameter, $0<\delta<1/2$, then
\begin{equation}\label{eq:sequential-information}
 \frac1M\sum_{j=1}^M\E_j\mathsf T
 \ge (1-\delta)\log_2M-h_2(\delta),
\end{equation}
where $h_2(x)=-x\log_2x-(1-x)\log_2(1-x)$. There is no deterministic
cap in this assertion. An infinite expected length satisfies it
trivially.
\end{lemma}
\begin{proof}
Give the parameter $V$ the uniform distribution and denote the
independent controller seed by $U$. Suppose the average expected length
is finite, so the terminal binary word $Z$ is finite almost surely.
For each fixed $U=u$, terminal words are prefix-free: after the rule
stops at one word it cannot also stop at a proper extension. Thus
$\sum_z2^{-|z|}\le1$. If this sum is $K_u>0$, comparison with the
probability distribution $2^{-|z|}/K_u$ by nonnegativity of relative
entropy gives
\[
 H(Z\mid U=u)\le\E(|Z|\mid U=u)+\log_2K_u
                    \le\E(\mathsf T\mid U=u).
\]
The countable case follows from the same log-sum inequality by
truncation. Integrate in $u$. Since $U$ is independent of $V$,
\[
 I(V;Z,U)=I(V;Z\mid U)\le H(Z\mid U)\le\E\mathsf T.
\]
Decode the parameter from the unique acceptable output set, using an
arbitrary index on failure. Its average error $p_e$ is at most $\delta$.
For the error indicator $E$, the elementary conditional-entropy bound is
\[
 H(V\mid Z,U)\le H(E\mid Z,U)+p_e\log_2(M-1)
 \le h_2(\delta)+\delta\log_2M.
\]
Subtract from $H(V)=\log_2M$ and use the preceding information bound.
The output randomness may be included in $U$ at the outset. Neither
continuous input values nor the word length yield an extra channel:
they are already determined by the seed and stopped word. The coding
and entropy principles used here are classical \cite{Shannon}; the
argument is included to specify the stopping-time scope exactly.
\end{proof}

\begin{theorem}[Expected attempted-bit lower bound]\label{thm:expected}
On the fixed physical class of Lemma~\ref{lem:packing}, any admissible
sequential procedure with success probability at least $3/4$ uniformly
in the centered loss \eqref{eq:accuracy-sets} satisfies
\begin{equation}\label{eq:expected-lower}
                 \sup_{\mathcal O}\E_{\mathcal O}\mathsf T
                         \ge c\nu^{-1/(s-2)}.
\end{equation}
The bound holds with exact implementation and arbitrary-precision
adaptive inputs, even with the fixed lattice, species and laboratory
frame disclosed. It therefore also holds for the stronger physical
reconstruction loss.
\end{theorem}
\begin{proof}
Choose $h\asymp\nu^{1/(s-2)}$ so that Lemma~\ref{lem:centered-packing}
gives separation greater than $2\nu$. Apply Lemma~\ref{lem:sequential}
with $M=2^m$ and $m\ge c/h$. The supremum dominates the average.
\end{proof}

At fixed confidence, the least uniform expected number of attempts on
any fixed class containing this packing lies between
$c\nu^{-1/(s-2)}$ and
$C\nu^{-1/(s-2)}\log^2(C/\nu)$, since the deterministic upper bound
is also an expected upper bound. The analogous deterministic-cap
statement remains \eqref{eq:matched-power}. These are two explicitly
proved resource criteria, not an identification of random length with
a deterministic cap. No sharp logarithmic or confidence factor follows.

\section{A resolution converse for table-dependent bounded errors}
\label{sec:calibration-floor}
The sufficient calibration scale in \eqref{eq:precision} has a converse
in the same bounded-error model. The converse uses physical start
perturbations, not corrupted labels added to an abstract membership
oracle. It concerns short collision commands; arbitrarily long flights
or additional observations of the actual launch position are not included.

\subsection{A common response produced by nearby physical tables}
For a table $\mathcal O_j$ denote its attempted hit bit by $B^j_a(q)$.
An admissible position-error implementation is a measurable map
$q\mapsto F_j(q,a)$ with $|F_j(q,a)-q|\le r$; time and direction
can remain exact. Such maps may depend on the table and the nominal
command. The controller sees only the resulting bit, not the actual
start. This is a submodel of the conditionally bounded, not necessarily
mean-zero couplings in Section~\ref{sec:setting}.

\begin{lemma}[Pointwise common response under two implementations]\label{lem:common-response}
Suppose $\mathcal O_0,\mathcal O_1$ consist of matched strictly convex
components of diameter at most $D_0$, with inradius at least $r_*>0$,
and that matched bodies have Hausdorff distance at most $e>0$.
In each table distinct bodies are separated by at least $d_0$.
Fix $t_{\max}<d_0$. If $e$ is sufficiently small compared with
$r_*$ and $d_0-t_{\max}$, there are measurable position perturbations
of size at most $2e$ such that, for every $|a|\le t_{\max}$ and
with closed-solid and swept-set conventions,
\begin{equation}\label{eq:common-response}
 B^0_a(F_0(q,a))=B^1_a(F_1(q,a))
                    =\min\{B^0_a(q),B^1_a(q)\}.
\end{equation}
The closed convention assigns zero on a solid boundary and one on a
swept boundary outside the solid. In every disagreement case the
forced-zero position is moved strictly into a solid or strictly outside
all sweeps before its response is evaluated. Agreement cases keep the
nominal point and use the same closed convention. The conclusion
applies to locally finite periodic arrays and to both the forward and
translated reverse commands.
\end{lemma}
\begin{proof}
For one fixed $a$ put $S_j(C)=C+[-a,0]$. Matched swept bodies have
Hausdorff distance at most $e$, because their support functions differ
by the same amount as those of the matched bodies. In either table,
\begin{equation}\label{eq:swept-separation}
 \dist(S_j(C),S_j(D))\ge d_0-|a|\quad(C\ne D).
\end{equation}
Indeed two points in the sweeps are $c-ua$ and $d-va$, with
$u,v\in[0,1]$. Their distance is at least $|c-d|-|u-v||a|$.
Thus a start belonging to a swept body determines that body uniquely.
Moreover no solid other than that body meets this sweep. On it the
hit bit is the indicator of the sweep minus the indicator of its body.

Leave the start unchanged when the two bits agree, and also on the
side whose bit is already zero. Consider $B^0_a(q)=1$, $B^1_a(q)=0$,
and let $C_0$ be the unique body whose sweep contains $q$. Write $C_1$
for its matched body. The separation \eqref{eq:swept-separation} and
Hausdorff matching imply the following exhaustive alternatives:
\[
                   q\in C_1,\qquad\hbox{or}\qquad q\notin\operatorname{int}S_1(C_1).
\]
To check this, every other sweep in table 1 is at distance at least
$d_0-|a|-e$ from $S_0(C_0)$, so cannot supply either a competing hit
or a solid start there. If $q\in S_1(C_1)\setminus C_1$, its bit
would be one.

In the first case $\dist(q,C_0)\le e$. Project $q$ onto $C_0$, and
move the projection a further distance less than $e/2$ toward a fixed
interior point of $C_0$. This produces an interior solid start within
$3e/2$ of $q$, whose bit is zero. Uniform inballs permit a fixed small
convex-combination factor. The construction is measurable.

In the second case, when $q$ is outside the closed sweep, let $q_1$
be its metric projection onto $S_1(C_1)$ and set
$v=(q-q_1)/|q-q_1|$. At a boundary point choose an outward unit normal
from its compact normal cone, measurably, for example by the least
angle in a fixed angular convention. In both cases
$q\cdot v\ge p_{S_1(C_1)}(v)$, and support matching gives
$p_{S_0(C_0)}(v)\le p_{S_1(C_1)}(v)+e$. Hence
\[
                       q'=q+\tfrac32 e v
\]
lies strictly outside $S_0(C_0)$. It cannot enter another swept body
of table 0, by \eqref{eq:swept-separation} and smallness of $e$.
It is therefore a free miss, with bit zero. The displacement is $3e/2$.
Projection onto a closed convex set is continuous, so this map is
measurable on its case domain. Reverse the indices when the differing
bit is on side 1. Local finiteness and the separation give measurable
component selection. For $a=0$ every bit is zero and no change is needed.
Boundary cases can be assigned by the same inward or outward moves;
in particular the final forced-zero positions are not on boundaries.

The proof is pointwise in the nominal start and displacement. Apply it
to $(q,a)$ for a forward command and to $(q+a,-a)$ for its reverse.
Fresh nominal draws and deterministic perturbation maps preserve
conditional independence. They do not supply another observable.
\end{proof}

\begin{theorem}[Minimax resolution under table-dependent bounded errors]
\label{thm:calibration-floor}
For each fixed $s=6+\beta$ there is a fixed physical subclass of
$\mathcal B_\eta$, with known primitive lattice, two known species
and the fixed laboratory frame, and constants $c,r_0>0$ such that
for every $0<r<r_0$ two tables in this class have centered-support
$C^2$ distance at least
\begin{equation}\label{eq:calibration-separation}
                            c r^{(s-2)/s}
\end{equation}
yet admit position-error implementations bounded by $r$ with identical
laws for the entire adaptive bit experiment. This holds for every
nominal command design with $|a|\le t_{\max}<d_0$, arbitrary input
precision and arbitrary numbers of attempts. A sequential stopping
rule based on the bits cannot distinguish them either.

To specify the uncertainty quantifiers, let $\mathfrak I_r(\mathcal O)$
be the measurable implementation maps permitted above, and let
$\mathcal A_{t_{\max}}$ be the admissible bit-based controllers. The
construction has the stronger, controller-uniform form
\begin{gather}\label{eq:calibration-quantifiers}
 \forall\,0<r<r_0\quad
 \exists\,(\mathcal O_0,F_0),(\mathcal O_1,F_1),\quad
             F_j\in\mathfrak I_r(\mathcal O_j),\notag\\
 \forall\,\mathcal A\in\mathcal A_{t_{\max}}:\qquad
 \mathsf{Law}_{\mathcal O_0,F_0}^{\mathcal A}(Z)
       =\mathsf{Law}_{\mathcal O_1,F_1}^{\mathcal A}(Z).
\end{gather}
Here $Z$ includes the returned bits and the controller's recorded
commands and stopping decisions, not the actual launch positions.
The maps $F_0,F_1$ need not agree and may use the nominal draw and
command. They are fixed measurable maps for the two worlds, not an
extra observation or a common known jitter distribution.

Consequently, uniform success probability greater than $1/2$ at
centered $C^2$ accuracy $\nu$ under all such position errors requires
\begin{equation}\label{eq:calibration-necessary}
                            r\le C\nu^{s/(s-2)}.
\end{equation}
Constants, not an equality threshold, are asserted. The result is a
worst-case bounded-calibration statement; it does not assert this
obstruction for one fixed common offset or for a known mean-zero law.
\end{theorem}
\begin{proof}
Use the two-species periodic construction in Lemma~\ref{lem:packing},
but compare the unit-disk support $p_0=1$ with a single nonzero bump
\[
 p_1(\theta)=1+a h^s\psi((\theta-\theta_0)/h).
\]
The amplitude $a$ is fixed small enough that all curvature, separation
and patch margins remain uniform. Choose
$h=(r/(8a\|\psi\|_\infty))^{1/s}$. Then the Hausdorff distance of
the two variable bodies is at most $e=r/8$. The fixed disk and all
periodic copies are matched without alteration. At the bump center
the second derivatives differ by $a|\psi''(0)|h^{s-2}$; the
first-harmonic estimate in Lemma~\ref{lem:centered-packing} shows that
centering removes only $O(h^s)$. This gives
\eqref{eq:calibration-separation}. The lattice and disclosed data are
identical for both tables; the bump does not change primitiveness.

Lemma~\ref{lem:common-response} produces two physical implementations
with position errors at most $2e<r$. Couple the parameter-independent
controller randomness and every nominal draw across the two worlds.
Their first bits are equal by \eqref{eq:common-response}. Inductively,
identical histories generate identical next commands and identical
bits. Thus the complete bit streams, nominal controller records and
all bit-measurable stopping decisions have identical laws. No finite
or infinite sequence of those observations distinguishes the worlds.
An output law common to two disjoint accuracy sets cannot place
probability greater than $1/2$ in each. If $2\nu$ is smaller than
\eqref{eq:calibration-separation}, uniform success greater than $1/2$
is impossible. Rearranging proves \eqref{eq:calibration-necessary}.
\end{proof}

\subsection{The two resource requirements}
Together with Theorems~\ref{thm:adaptive} and \ref{thm:expected}, this
gives two necessary resources on one fixed nondegenerate subclass:
\[
 \sup_{\mathcal O}\E_{\mathcal O}\mathsf T
       \ge c\nu^{-1/(s-2)},\qquad
 r\le C\nu^{s/(s-2)}.
\]
Conversely the reciprocal compass construction attains the geometric
accuracy with the stated logarithmic overhead in attempts when
$\ell+\tau+t\alpha\le c\nu^{s/(s-2)}$. The converse calibration
necessity already holds with duration and angle exact, so it concerns
position error alone, not separate optimal exponents for all three
actuators. The lower bound allows any directions of length at most
$t_{\max}$, a larger design class than the four-atom upper construction.
The finite digital realization in Section~\ref{sec:digital} does not
alter either bound.

Thus the shrinking physical tolerance is not merely an omitted term in
a sample count: under the stipulated worst-case coupling it cannot be
replaced by additional samples. This is a statement within the active
collision model. Periodicity, the bounded presentation and the known
patch margin are still priors for the uniform global conclusion.
Neither crystallinity, apparatus calibration nor passive spectral
rigidity is inferred from the returned bits.

\subsection*{Source history and assisted analysis}
The article contains all proofs required for its statements. Earlier
localized, finite-control and calibration results are retained in the
appendices; the reviewed source packages remain at their original
repository paths. They are not independently published articles.
AI-assisted analysis contributed to the collision-query, footprint and
scale-recovery arguments and to the source checks. The named author is
responsible for the proofs and references. Finite diagnostics and a
successful build are reproducibility evidence and do not certify
continuum proofs or validate a physical apparatus.
\begin{thebibliography}{99}
\raggedright
\bibitem{CN}
R. M. Castro and R. D. Nowak, \emph{Minimax bounds for active learning},
IEEE Trans. Inform. Theory \textbf{54} (2008), 2339--2353.
\href{https://doi.org/10.1109/TIT.2008.920189}{doi:10.1109/TIT.2008.920189}.
\bibitem{LCK}
A. Locatelli, A. Carpentier, and S. Kpotufe,
\emph{An adaptive strategy for active learning with smooth decision boundary},
Proceedings of Algorithmic Learning Theory, Proceedings of Machine
Learning Research \textbf{83} (2018), 547--571.
\bibitem{LawlerLimic}
G. F. Lawler and V. Limic, \emph{Random Walk: A Modern Introduction},
Cambridge Studies in Advanced Mathematics \textbf{123}, Cambridge
University Press, 2010.
\bibitem{Shannon}
C. E. Shannon, \emph{A mathematical theory of communication},
Bell System Tech. J. \textbf{27} (1948), 379--423, 623--656.
\bibitem{Retained}
Q. Qi, \emph{Scalar collision laws and recognition of periodic dispersing
billiards}, A2 v31, October 4, 2026. Complete reviewed manuscript retained in the repository source archive;
not an independently published article or a required supplementary volume.
\bibitem{Schneider}
R. Schneider, \emph{Convex Bodies: The Brunn--Minkowski Theory},
second expanded edition, Encyclopedia of Mathematics and its Applications
\textbf{151}, Cambridge University Press, 2014.
\bibitem{BKY}
V.-E. Brunel, J. M. Klusowski, and D. Yang,
\emph{Estimation of convex supports from noisy measurements},
Bernoulli \textbf{27} (2021), 772--793.
\href{https://doi.org/10.3150/20-BEJ1229}{doi:10.3150/20-BEJ1229}.
\bibitem{YCJ}
S. Yan, K. Chaudhuri, and T. Javidi,
\emph{Active learning from imperfect labelers},
Advances in Neural Information Processing Systems \textbf{29} (2016).
\bibitem{GK}
P. W. Goldberg and S. Kwek,
\emph{The precision of query points as a resource for learning convex
polytopes with membership queries},
Proceedings of the Thirteenth Annual Conference on Computational
Learning Theory, Morgan Kaufmann, 2000, 225--235.
\end{thebibliography}

\end{document}
